\chapter{Arithmetic forms and operator domains}
\label{chap:weil-operator-domains}
\begingroup
\noeqref{eq:test-class,
eq:operations,
eq:centered-mellin,
eq:target,
eq:strip-decay,
eq:gaussian,
eq:inverse,
eq:star-transform,
eq:euler,
eq:gamma-product,
eq:reciprocal-growth,
eq:digamma,
eq:gamma-real,
eq:theta-reflection,
eq:theta-mellin,
eq:entire-completion,
eq:entire-growth,
eq:boundary-pole,
eq:boundary-product,
eq:strict-strip,
eq:log-reflection,
eq:horizontal-H,
eq:functional,
eq:spine-contour,
eq:prime-inversion,
eq:gaussian-explicit,
eq:gaussian-inverse,
eq:polynomial-transform,
eq:gaussian-strip,
eq:positivity,
eq:negative-pair,
eq:separator-tail,
eq:conv,
eq:transform,
eq:fourier,
eq:square-transform,
eq:gamma-normalization,
eq:mellin,
eq:weil-functional,
eq:explicit,
eq:hermitian,
eq:weil-criterion,
eq:zero-norm,
eq:gns-action,
eq:adjoint-inclusion,
eq:multiplier-bound,
eq:all-moments,
eq:common-domain,
eq:closed-products,
eq:graph-control,
eq:pullback,
eq:gns-map,
eq:completion-map,
eq:algebraic-intertwining,
eq:cstar-bounds,
eq:bounded-intertwining,
eq:graph-transfer,
eq:ambient-transfer,
eq:unit-extension,
eq:unit-criterion,
eq:minimal-mass,
eq:approx-vector,
eq:cyclic-vector,
eq:hilbert-bound,
eq:gaussian-semigroup,
eq:weil-gaussian-asymptotic,
eq:weil-unitization-divergence,
eq:prime-gaussian-bound,
eq:gaussian-gram-limit,
eq:limiting-gaussian-gram,
eq:indefinite-negative,
eq:factorization,
eq:weil-gaussian-moments}
\newcommand{\A}{\mathcal A}

Closing an operator kernel to a scalar and giving states positive norms
are separate requirements in a field theory.
A cyclic trace controls an operator contraction.  An involution and a
positive functional instead produce a Hilbert space and operators on a
common domain.  Arithmetic supplies a concrete test of the second
construction: can its prime, archimedean and polar terms define a positive
pairing without prescribing the spectrum in advance?

The arithmetic terms of the zeta explicit formula define a Hermitian
pairing on a nonunital convolution algebra.  Gaussian reflection and a
contour integral derive the full formula on this algebra.  A separate
polynomial--Gaussian argument identifies its positivity with the Riemann
hypothesis.  A local Gaussian calculation proves a further
restriction: the full Weil functional admits no positive unitization of
finite mass and no realization by a finite-norm vector on an invariant
$*$-representation domain.  The same calculation gives positive Gaussian
subspaces of every finite dimension at sufficiently small scales.
We construct the common graph completion of
nonunital GNS multipliers and prove its exact domain and transfer
properties.  Closed quadratic forms for individual multipliers need not
have invariant domains.  Their common domain over the whole algebra does.
An independent positive arithmetic functional remains open, with its
possible representation constrained by the Gaussian divergence.


The decisive question is whether the arithmetic expression is a squared
norm for a reason independent of the location of the zeros.
Such a norm would construct a Hilbert space from test functions.
It need not come from a distinguished vector representing the functional.
The Gaussian calculation in Section~\ref{sec:weil-mass} excludes that
stronger realization.  We first derive the arithmetic expression and its
positivity criterion; then we construct the representation and its domain
under a stated positivity hypothesis.

\section{The arithmetic identity and its variables}

The Dirichlet series $\sum_{n\geq1}n^{-s}$ has two useful descriptions.
Prime factorization gives a product when $\operatorname{Re}s>1$.
A Gaussian sum gives a continuation and a reflection in $s=1/2$.
Taking a logarithmic derivative changes prime factors into weighted prime
powers and changes zeros into simple poles.  A contour integral can then
compare them.  The test function must control that contour before it can
define an identity.

All vector spaces are complex, and all real integrals use Lebesgue measure.
The starting language is calculus, absolutely convergent integrals and
series, and the complete fields $\R$ and $\C$.
We use the usual limit rules for absolutely integrable functions: an
integrable common bound permits passage through an integral.
Every infinite sum and parameter integral below has an explicit such bound.
The elementary complex-variable tools needed for the contour are recorded
in Section~\ref{sec:complex}.

For $a>0$, define
\begin{equation}
 \begin{aligned}
 \A_a&=\left\{g\in C^\infty(\R):
   p_{a,k}(g):=\int_{\R}e^{a|u|}|g^{(k)}(u)|\,du<\infty
                     \text{ for every }k\geq0\right\},\\
 \A&=\bigcup_{a>1/2}\A_a.
 \end{aligned}
 \label{eq:test-class}
\end{equation}
Here $C^\infty$ means that every real derivative exists and is continuous.
The functions may be complex-valued and need not be even.
Define convolution, involution, and Fourier transform by
\begin{equation}
 (f\star g)(u)=\int_{\R}f(v)g(u-v)\,dv,
 \qquad g^*(u)=\overline{g(-u)},
 \qquad \widehat g(z)=\int_{\R}g(u)e^{izu}\,du.
 \label{eq:operations}
\end{equation}
The complex variable suited to the Dirichlet series is
\begin{equation}
 F_g(s)=\widehat g\bigl(-i(s-1/2)\bigr)
       =\int_{\R}g(u)e^{(s-1/2)u}\,du
       =\int_0^\infty g(\log x)x^{s-1/2}\,\frac{dx}{x}.
 \label{eq:centered-mellin}
\end{equation}
The last equality is the change of variable $x=e^u$.
In particular $F_g(1/2+it)=\widehat g(t)$.
The half in this substitution determines the weights of all prime terms.

Write $\Lambda_{\mathrm{vM}}(n)=\log p$ if $n=p^m$ with $p$ prime and $m\geq1$,
and write $\Lambda_{\mathrm{vM}}(n)=0$ otherwise.
The gamma function and the meromorphic continuation of $\zeta$ will be
constructed below.  Put $\psi=\Gamma'/\Gamma$.
The identity to be proved is
\begin{equation}
\begin{split}
 \sum_\rho F_g(\rho)={}&F_g(0)+F_g(1)\\
 &+\frac1{2\pi}\int_{\R}\widehat g(t)
 \left[\operatorname{Re}\psi\left(\frac14+\frac{it}{2}\right)-\log\pi\right]dt\\
 &-\sum_{n\geq2}\frac{\Lambda_{\mathrm{vM}}(n)}{\sqrt n}
                      \bigl(g(\log n)+g(-\log n)\bigr).
\end{split}
\label{eq:target}
\end{equation}
Zeros $\rho$ are counted with multiplicity.  Their precise definition will
exclude the canceled gamma poles and the poles at zero and one.
Theorem~\ref{thm:explicit} proves \eqref{eq:target} for every $g\in\A$.
The proof uses no positivity statement.

\section{Elementary complex estimates}
\label{sec:complex}

A holomorphic function is a complex-differentiable function on an open set.
Differentiability implies continuity, but we will not assume continuity
of the derivative in deriving the integral formulas.
A meromorphic function is locally a quotient of holomorphic functions,
with denominator not identically zero.
A zero or pole has order $m$ if its local expression contains a factor
$(z-z_0)^m$ or $(z-z_0)^{-m}$ with nonzero remaining factor.
The residue is the coefficient of $(z-z_0)^{-1}$.

For a piecewise continuously differentiable path $\gamma:[a,b]\to\C$,
define
\[
 \int_\gamma h(w)\,dw
   =\int_a^b h(\gamma(t))\gamma'(t)\,dt,
\]
adding the integrals over its differentiable pieces.
Reversing a path reverses the sign of its integral.

\begin{lemma}[Integration around a triangle]
\label{lem:weil-triangle-integral}
If $h$ is holomorphic on a neighborhood of a closed triangle $T$, then
$\int_{\partial T}h(w)\,dw=0$, with the boundary oriented counterclockwise.
\end{lemma}
\begin{proof}
Joining the edge midpoints divides $T$ into four triangles.
Their oriented boundary integrals sum to the original integral because
every internal edge occurs in both directions.
Choose a subtriangle whose integral has modulus at least one quarter
of that of its parent, and repeat.  The resulting nested closed triangles
$T_n$ have diameter $2^{-n}D$ and perimeter $2^{-n}P$, where $D,P$ are
those of $T$.  Points chosen in these triangles form a Cauchy sequence.
Completeness gives a limit $w_0$ belonging to every $T_n$, and their
shrinking diameters make this common point unique.

Differentiability at $w_0$ gives
\[
 h(w)=h(w_0)+h'(w_0)(w-w_0)+(w-w_0)\epsilon(w),
 \qquad \epsilon(w)\longrightarrow0\quad(w\to w_0).
\]
The first two terms have polynomial primitives, so their integrals around
each triangle vanish by the real fundamental theorem of calculus on each
edge.  Put $\epsilon_n=\sup_{w\in T_n}|\epsilon(w)|$, defining
$\epsilon(w_0)=0$.  Then $\epsilon_n\to0$, and
\[
 \left|\int_{\partial T}h(w)\,dw\right|
 \leq4^n\left|\int_{\partial T_n}h(w)\,dw\right|
 \leq4^n\epsilon_n(2^{-n}D)(2^{-n}P)
 =\epsilon_nDP\longrightarrow0.
\]
\end{proof}

\Needspace{5\baselineskip}
\begin{lemma}[Contour formulas]
\label{lem:cauchy}
Let a function be holomorphic on a neighborhood of a closed disc.
For any interior point $z$,
\[
 h(z)=\frac1{2\pi i}\int_{|w-z_0|=R}\frac{h(w)}{w-z}\,dw,
 \qquad
 |h^{(k)}(z_0)|\leq \frac{k!}{R^k}\max_{|w-z_0|=R}|h(w)|.
\]
The contour has counterclockwise orientation.
For a meromorphic function on a rectangle with no boundary poles, its
boundary integral is $2\pi i$ times the sum of its interior residues.
Locally uniform limits of holomorphic functions are holomorphic.
\end{lemma}

\begin{proof}
Summing the triangle identity gives zero around any polygonal region
on whose neighborhood the integrand is holomorphic, including a region
with polygonal holes.  To obtain a finite subdivision, draw parallel
lines through its finitely many vertices, divide the resulting strips
into trapezoids and triangles, and split each trapezoid along a diagonal.
The internal edges cancel, leaving the outer boundary counterclockwise
and the inner boundaries clockwise.

The same identity holds for a circular boundary by polygonal
approximation.  Linear interpolation of its smooth parametrization
converges uniformly in position; on each interpolation interval its
derivative is the average of the original derivative.  Uniform continuity
of that derivative gives convergence of tangents in the integral norm.
Continuity of the integrand, uniformly on a compact neighborhood of the
circle, then gives convergence of the boundary integrals.

Apply this identity to $h(w)/(w-z)$ between the original circle and a
small circle centered at $z$.  Its neighborhood contains no singularity
of this integrand.  On the small circle, write
$h(w)=h(z)+O(|w-z|)$ and let its radius tend to zero.
The constant term integrates to $2\pi i h(z)$ and the error tends to zero.
This proves the first formula without a regularity assumption on $h'$.
Expansion of $(w-z)^{-1}$ as a geometric series gives the derivative formula
and the displayed estimate, as well as a convergent local power series
for every holomorphic function.
For the residue formula, remove disjoint small discs around the poles.
The same boundary argument applies to the remaining region.
The finite principal part at each pole follows by factoring the first
nonzero Taylor term of the denominator in a local quotient.
On its small circle, only the term $(w-w_0)^{-1}$ has nonzero integral,
namely $2\pi i$ times its coefficient.  Adding the circles proves the
residue statement.
For a locally uniform limit, pass to the limit in the circle formula.
The resulting expression is differentiable in the interior point and gives
its convergent power series.  These arguments also justify the formulas
on discs and rectangles with finitely many small discs removed.
\end{proof}

The circle formula implies the identity theorem: a nonzero holomorphic
function has isolated zeros, since its first nonzero Taylor coefficient
factors out the entire order of a zero.  It also implies the maximum
modulus principle.  If an interior modulus is maximal, equality in the
circle average forces constant argument and modulus on a small circle.
Its Taylor coefficients of positive degree then vanish, and the function
is constant on the connected domain.

\Needspace{5\baselineskip}
\begin{lemma}[Logarithms and counting zeros]
\label{lem:jensen}
A zero-free holomorphic function on a disc has a holomorphic logarithm.
If an entire function $h$ satisfies $h(0)\ne0$ and
\[
 \log\max_{|z|\leq R}|h(z)|\leq C R\log(R+2)\qquad(R\geq2),
\]
then the number $n_h(R)$ of its zeros in $|z|\leq R$, with multiplicity,
satisfies $n_h(R)=O(R\log(R+2))$.
\end{lemma}

\begin{proof}
For a zero-free function, the integral of $h'/h$ around any closed contour
in the disc is zero by Lemma~\ref{lem:cauchy}.  Its path integral is thus a
primitive $\ell$.  Differentiation shows that $e^{-\ell}h$ is constant.
Choosing the additive constant gives $h=e^\ell$.

Choose $r>R$ so that $h$ has no zero on $|z|=r$.
Factor out its finitely many zeros $a_j$ inside that circle.
The remaining nonzero factor has a logarithm whose real part has its value
at zero equal to its circle average.
For $|a|<r$, the power series of $\log(1-ae^{-i\theta}/r)$ gives
the circle average $\log|re^{i\theta}-a|=\log r$.
Subtracting the value at zero gives
\[
 \sum_{|a_j|<r}\log\frac r{|a_j|}
 =\frac1{2\pi}\int_0^{2\pi}\log|h(re^{i\theta})|\,d\theta
       -\log|h(0)|.
\]
Take $2R<r<2R+1$, avoiding the isolated radii of zeros.
Each zero in $|z|\leq R$ contributes at least $\log2$.
The asserted growth bound proves the count.
\end{proof}

The contour will need a logarithmic derivative estimate, not only a count.
The next lemma supplies that extra information locally.

\Needspace{5\baselineskip}
\begin{lemma}[A local logarithmic derivative bound]
\label{lem:local-log}
Let $h$ be holomorphic on a neighborhood of $|z|\leq R$, where
$5\leq R\leq6$, and suppose $h(0)\ne0$ and $h$ has no boundary zero.
Suppose
\[
 |h(z)|\leq e^M\quad(|z|\leq R),\qquad |h(0)|\geq e^{-M},\qquad M\geq1.
\]
If $h$ has $N$ zeros in the disc and $|z|\leq3$ is at distance at least
$\delta>0$ from each zero, then
\[
 \left|\frac{h'(z)}{h(z)}\right|
 \leq C M+N(\delta^{-1}+1/2),
\]
where $C$ is an absolute constant.
\end{lemma}

\begin{proof}
For each zero $a$, including multiplicities, use the factor
$b_a(z)=R(z-a)/(R^2-\bar a z)$.
It has modulus one on $|z|=R$ and no denominator zero in the disc.
Divide $h$ by their product and remove the canceled singularities.
The quotient $v$ is holomorphic and zero-free.
The maximum principle gives $|v|\leq e^M$.
Since $|b_a(0)|<1$, also $|v(0)|\geq|h(0)|\geq e^{-M}$.
Choose $\ell=\log v$ with $\ell(0)$ real after multiplying $v$ by a constant
of modulus one.  The function $p=M+1-\ell$ has positive real part and
$1\leq p(0)\leq2M+1$.

The function $q=(p-p(0))/(p+p(0))$ maps the disc into the unit disc and
vanishes at zero.  Applying the maximum principle to $q(z)/z$ on smaller
circles and increasing their radii to $R$ gives $|q(z)|\leq|z|/R$.
Hence $|p(z)|\leq9p(0)$ on $|z|\leq4$.
The derivative estimate on radius-one circles about $|z|\leq3$ yields
$|\ell'(z)|=|p'(z)|\leq9p(0)\leq27M$.
Finally
\[
 \frac{h'}h=\ell'+\sum_a
 \left(\frac1{z-a}+\frac{\bar a}{R^2-\bar a z}\right).
\]
For $|z|\leq3$, the second fraction has modulus at most $1/(R-3)\leq1/2$.
The assumed distance bound controls the first.
\end{proof}

\Needspace{7\baselineskip}
\section{The test algebra and Gaussian Fourier inversion}
\label{sec:tests}

\begin{lemma}[Decay, boundary values, and transforms]
\label{lem:test-bounds}
If $g\in\A_a$, then $e^{a|u|}g^{(k)}(u)\to0$ at both ends of $\R$,
and $|g^{(k)}(u)|\leq C_k e^{-a|u|}$.
The transform $\widehat g$ is holomorphic for $|\operatorname{Im}z|<a$.
For every $0<b<a$ and every integer $N\geq0$,
\begin{equation}
 \sup_{|\operatorname{Im}z|\leq b}(1+|z|)^N|\widehat g(z)|<\infty.
 \label{eq:strip-decay}
\end{equation}
Equivalently, $F_g$ is holomorphic for $1/2-a<\operatorname{Re}s<1/2+a$
and decreases faster than any inverse power of $|\operatorname{Im}s|$
on every smaller closed vertical strip.
\end{lemma}

\begin{proof}
On $u\geq0$, the function $h(u)=e^{au}g^{(k)}(u)$ and its derivative
$e^{au}(ag^{(k)}+g^{(k+1)})$ are integrable.
The integral of the derivative makes $h$ have a finite limit at infinity.
Integrability of $h$ forces that limit to be zero.
Moreover $|h(u)|\leq\int_u^\infty|h'(v)|\,dv$.
The negative half-line has the same proof after reflection.

For $b<a$, the functions $|u|^m e^{b|u|}|g(u)|$ are integrable for every
$m$, because $|u|^m e^{-(a-b)|u|}$ is bounded.
These bounds permit termwise complex differentiation of the integral.
Repeated integration by parts has zero boundary terms by the first assertion
and gives
\[
 (-iz)^N\widehat g(z)=\int_{\R}g^{(N)}(u)e^{izu}\,du.
\]
Its right side is bounded on $|\operatorname{Im}z|\leq b$.
Combining this with the bound for $N=0$ proves \eqref{eq:strip-decay}.
Substitution in \eqref{eq:centered-mellin} proves the last assertion.
\end{proof}

\Needspace{5\baselineskip}
\begin{lemma}[The normalized Gaussian]
\label{lem:gaussian}
For $A>0$, put
\[
 g_A(u)=(4\pi A)^{-1/2}e^{-u^2/(4A)}.
\]
Then $g_A\in\A_a$ for every $a>0$, and
\begin{equation}
 \widehat g_A(z)=e^{-Az^2},\qquad
 g_A\star g_B=g_{A+B},\qquad \int_{\R}g_A(u)\,du=1.
 \label{eq:gaussian}
\end{equation}
For $g\in\A$ one has the inverse formula
\begin{equation}
 g(u)=\frac1{2\pi}\int_{\R}\widehat g(t)e^{-itu}\,dt.
 \label{eq:inverse}
\end{equation}
The same inverse formula holds for a smooth function $g$ with
$g,g',g''\in L^1(\R)$.
\end{lemma}

\begin{proof}
Squaring $\int e^{-x^2}\,dx$ and using polar coordinates gives its value
$\sqrt\pi$.  Scaling proves the mass assertion.
Every derivative of $g_A$ is a polynomial times this Gaussian, proving
membership and vanishing of every exponential-weighted boundary value.
For real $t$, differentiate $J(t)=\widehat g_A(t)$ under the integral.
Since $u g_A=-2A g_A'$, integration by parts gives
$J'(t)=-2AtJ(t)$.  The initial value $J(0)=1$ gives $J(t)=e^{-At^2}$.
Both sides are entire, so the identity theorem extends the formula to $z\in\C$.
Completing the square in the convolution integral gives $g_{A+B}$.

For inversion, insert the factor $e^{-At^2}$ on the right of
\eqref{eq:inverse}.  Absolute integrability permits reversing the integrals.
The Gaussian transform just proved gives
\[
 \frac1{2\pi}\int_{\R}e^{-At^2}\widehat g(t)e^{-itu}\,dt
       =(g_A\star g)(u).
\]
The left side tends to the proposed inverse by \eqref{eq:strip-decay}.
On the right, continuity of $g$ controls $|v|<\eta$ in the convolution
difference, and the Gaussian mass outside $|v|<\eta$ tends to zero.
Boundedness of $g$ controls that tail.  Thus $(g_A\star g)(u)\to g(u)$.
For the stated larger class, integrability of $g,g',g''$ makes $g,g'$
bounded with zero limits at infinity, by the half-line argument of
Lemma~\ref{lem:test-bounds} with zero weight.
Two integrations by parts give an integrable bound for $\widehat g$.
The same proof then applies.
\end{proof}

The parameter $A$ is also a heat parameter.  Direct differentiation gives
\[
 \partial_A g_A(u)=\left(-\frac1{2A}+\frac{u^2}{4A^2}\right)g_A(u)
                  =\partial_u^2g_A(u).
\]
Thus convolution by $g_A$ solves heat evolution on the real line for the
smooth integrable initial functions considered here.
The factor $e^{-At^2}$ damps high Fourier frequencies, and
$g_A\star g_B=g_{A+B}$ expresses addition of elapsed heat times.
This supplies an explicit analytic regularization for the inverse transform
and for the arithmetic identities below.

A complex $*$-algebra is an associative complex algebra with a conjugate-linear
map $a\mapsto a^*$ such that $(a^*)^*=a$ and $(ab)^*=b^*a^*$.
For the $*$-algebras in this chapter, a unit is a nonzero element $1$
with $1a=a1=a$ for every $a$;
nonunital means that no such element is present.
In particular, the zero algebra is nonunital under this convention.

\Needspace{5\baselineskip}
\begin{proposition}[Algebra and change of exponential convention]
\label{prop:algebra}
Each $\A_a$ for $a>0$, and also $\A$, is a commutative nonunital
$*$-algebra under \eqref{eq:operations}, and
\begin{equation}
 \widehat{f\star g}(z)=\widehat f(z)\widehat g(z),\qquad
 \widehat{f^*}(z)=\overline{\widehat f(\bar z)}.
 \label{eq:star-transform}
\end{equation}
These equations hold on every common strip of holomorphy.
The union in \eqref{eq:test-class} is unchanged if weighted $L^1$ bounds
are replaced by weighted supremum bounds, with the exponential rate allowed
to decrease strictly.
\end{proposition}

\begin{proof}
For $a>0$ and $f,g\in\A_a$, use
$e^{a|u|}\leq e^{a|v|}e^{a|u-v|}$ to obtain
$p_{a,k}(f\star g)\leq p_{a,k}(f)p_{a,0}(g)$.
If $f,g\in\A$, they admit a common $a>1/2$.
Differentiation, reassociation, and exchange of factors are justified by
these absolute integral bounds.  Reflection proves the involution rule,
and reversing the absolutely convergent double integral proves the first
transform equation.  The second follows from $u\mapsto-u$ and conjugation.
If an identity $e$ existed in any one of these algebras, convolution with
$g_A$ would imply $\widehat e(t)=1$, contradicting \eqref{eq:strip-decay}.

Lemma~\ref{lem:test-bounds} changes weighted $L^1$ bounds into weighted
supremum bounds at the same rate.  Conversely, if
$|g^{(k)}(u)|\leq C_k e^{-a|u|}$, then for $1/2<b<a$,
$p_{b,k}(g)\leq2C_k/(a-b)$.
This proves the union equality without identifying its individual fixed-rate
spaces or silently removing the loss in the rate.
\end{proof}

\section{Prime factors and the gamma factor}
\label{sec:euler-gamma}

For positive $x$ and complex $s$, write $x^s=\exp(s\log x)$ with the real
logarithm of $x$.  This convention fixes every power used below.

\Needspace{5\baselineskip}
\begin{lemma}[The convergent prime product]
\label{lem:euler}
The series $\zeta(s)=\sum_{n\geq1}n^{-s}$ defines a holomorphic nonzero
function for $\operatorname{Re}s>1$.  In that half-plane,
\begin{equation}
 \zeta(s)=\prod_p(1-p^{-s})^{-1},\qquad
 \frac{\zeta'}{\zeta}(s)=-\sum_{n\geq2}\Lambda_{\mathrm{vM}}(n)n^{-s}.
 \label{eq:euler}
\end{equation}
Both the logarithmic series and its derivative converge absolutely on every
smaller closed half-plane.  On $s=2+it$, $|\zeta(s)|\geq\zeta(2)^{-1}$.
\end{lemma}

\begin{proof}
Existence and uniqueness of prime factorization follow by induction from
the fact that a prime dividing a product divides a factor.  The latter
follows from the integer Euclidean algorithm and its identity for the
greatest common divisor.
Expanding a finite prime product into geometric series therefore gives the
sum over integers supported on those primes.
For $\sigma>1$, the integral comparison
$\sum n^{-\sigma}\leq1+\int_1^\infty x^{-\sigma}\,dx$ is finite.
It bounds the omitted tails and proves the product identity.

The series
$\sum_p\sum_{m\geq1}p^{-ms}/m$ and its derivative are locally uniformly
absolutely convergent: their absolute values are bounded by constants times
$\sum_{n\geq2}(1+\log n)n^{-\sigma}$.
Exponentiating that series gives the product, so the product is nonzero.
Differentiation gives \eqref{eq:euler}.
Finally
\[
 |\zeta(2+it)|^{-1}
 \leq\prod_p(1+p^{-2})\leq\prod_p(1-p^{-2})^{-1}=\zeta(2).
\]
\end{proof}

\Needspace{5\baselineskip}
\begin{lemma}[Gamma product and vertical bounds]
\label{lem:gamma}
For $\operatorname{Re}z>0$, the integral
\[
 \Gamma(z)=\int_0^\infty e^{-x}x^{z-1}\,dx
\]
is holomorphic and nonzero.  Its reciprocal extends to the entire function
\begin{equation}
 E(z)=\frac1{\Gamma(z)}
 =z e^{\gamma z}\prod_{n\geq1}(1+z/n)e^{-z/n},
 \label{eq:gamma-product}
\end{equation}
where $\gamma=\lim_{N\to\infty}(\sum_{n=1}^N1/n-\log N)$.
Its zeros are simple and are exactly $0,-1,-2,\ldots$.
For $R\geq2$,
\begin{equation}
 |E(z)|\leq\exp(CR\log(R+2))\quad(|z|\leq R).
 \label{eq:reciprocal-growth}
\end{equation}
The logarithmic derivative satisfies
\begin{equation}
 \psi(z)=-\gamma+\sum_{n\geq0}
           \left(\frac1{n+1}-\frac1{n+z}\right),
 \qquad
 \psi(\sigma+it)=O_{\alpha,\beta}(\log(2+|t|))
 \label{eq:digamma}
\end{equation}
uniformly for $0<\alpha\leq\sigma\leq\beta$.
\end{lemma}

\begin{proof}
The integral and its derivatives have integrable majorants on compact
subsets of $\operatorname{Re}z>0$.
Repeated real integration by parts gives
\[
 \int_0^1 x^{z-1}(1-x)^N\,dx
       =\frac{N!}{z(z+1)\cdots(z+N)}.
\]
After multiplying by $N^z$ and substituting $u=Nx$, the integrand tends to
$u^{z-1}e^{-u}$.  The inequality $(1-u/N)^N\leq e^{-u}$ for $0\leq u<N$
is an integrable bound, also after any fixed number of $z$ derivatives
on compact subsets.  Thus
\[
 \Gamma(z)=\lim_{N\to\infty}
        \frac{N^zN!}{z(z+1)\cdots(z+N)}.
\]
The sequence $\sum_{n=1}^N1/n-\log N$ decreases and is bounded below
by the comparison of $1/x$ with its left and right rectangle sums.
Taking reciprocals of the finite expressions gives the finite products
in \eqref{eq:gamma-product}.
For $|w|\leq1/2$, the power series of the logarithm gives
$|\log(1+w)-w|\leq C|w|^2$.
Hence their tails converge uniformly on compact sets to a nonzero factor.
The finite initial factors identify the zeros and their orders.
Multiplication of the two limits gives $\Gamma E=1$ on the right half-plane,
proving nonvanishing and the reciprocal formula.

For $|z|\leq R$, split the product at $n=\lceil2R\rceil$.
The logarithms of absolute values of the initial factors are bounded above by
\[
 \sum_{n\leq\lceil2R\rceil}
       \bigl(\log(1+R/n)+R/n\bigr)=O(R\log(R+2)).
\]
The remaining logarithmic tail is $O(R^2\sum_{n>2R}n^{-2})=O(R)$.
Including $z e^{\gamma z}$ proves \eqref{eq:reciprocal-growth}.

Logarithmic differentiation gives the series in \eqref{eq:digamma}.
For $|t|\geq1$, split that series at $n=\lceil2(|t|+\beta+1)\rceil$.
The sum of $1/(n+1)$ before that point is $O(\log(2+|t|))$,
and the sum of $|n+z|^{-1}$ is $O(1)$ because $|n+z|\geq|t|$.
After that point, each difference is bounded by $C|z-1|/(n+1)^2$,
whose sum is $O(1)$.
The compact part with $|t|\leq1$ has no pole and is bounded.
\end{proof}

Define
\begin{equation}
 \Gamma_{\R}(s)=\pi^{-s/2}\Gamma(s/2),\qquad
 H_\Gamma(s)=\frac{\Gamma_{\R}'}{\Gamma_{\R}}(s)
          =-\frac12\log\pi+\frac12\psi(s/2).
 \label{eq:gamma-real}
\end{equation}
The half in $s/2$ and the real factor $\pi^{-s/2}$ will produce respectively
the argument $1/4+it/2$ and the term $-\log\pi$ in \eqref{eq:target}.

\section{Gaussian reflection and the completed function}
\label{sec:theta}

The completion of the Dirichlet series will come from a lattice sum.
The identity needed for that sum is a direct consequence of the Gaussian
Fourier transform, with its normalization kept explicit.

The normalization in Lemma~\ref{lem:gaussian} also gives the Fourier
integral used by periodization.  The substitution
$A=1/(4\pi x)$ and $z=-2\pi t$ relates the two conventions.

\Needspace{5\baselineskip}
\begin{lemma}[Gaussian Fourier transform]
\label{lem:gaussian-real-transform}
For $x>0$ and real $t$,
\[
 \int_{\R}e^{-\pi xv^2}e^{-2\pi itv}\,dv
       =x^{-1/2}e^{-\pi t^2/x}.
\]
\end{lemma}
\begin{proof}
Squaring $I=\int_{\R}e^{-u^2}\,du$ and using polar coordinates gives
$I^2=2\pi\int_0^\infty re^{-r^2}\,dr=\pi$.
The integrand is positive, so $I=\sqrt\pi$.
Scaling gives $\int_{\R}e^{-\pi xv^2}\,dv=x^{-1/2}$.
Let $J(t)$ denote the integral in the statement.
Every fixed power of $|v|$ times $e^{-\pi xv^2}$ is integrable.
This bound permits differentiation in $t$ and integration by parts in $v$.
Using $\partial_v e^{-\pi xv^2}=-2\pi xv e^{-\pi xv^2}$ gives
\[
 J'(t)=\frac{i}{x}\int_{\R}
       \partial_v(e^{-\pi xv^2})e^{-2\pi itv}\,dv
       =-\frac{2\pi t}{x}J(t).
\]
The boundary terms vanish because of Gaussian decay.
The initial value $J(0)=x^{-1/2}$ solves this differential equation
and proves the formula.
\end{proof}

\Needspace{5\baselineskip}
\begin{lemma}[Periodization of a Gaussian]
\label{lem:theta}
For $x>0$, put $\theta(x)=\sum_{n\in\Z}e^{-\pi n^2x}$.
Then
\begin{equation}
 \theta(x)=x^{-1/2}\theta(1/x).
 \label{eq:theta-reflection}
\end{equation}
For $x\geq1$, one has $0<\theta(x)-1\leq C e^{-\pi x}$.
\end{lemma}

\begin{proof}
The periodized function
$P_x(u)=\sum_{n\in\Z}e^{-\pi x(u+n)^2}$ and all its derivatives converge
uniformly for $0\leq u\leq1$.
Its $k$th Fourier coefficient, with period one, is
\[
 \int_0^1P_x(u)e^{-2\pi iku}\,du
 =\int_{\R}e^{-\pi xv^2}e^{-2\pi ikv}\,dv
 =x^{-1/2}e^{-\pi k^2/x},
\]
by Lemma~\ref{lem:gaussian}.
These coefficients are absolutely summable and define a continuous
periodic Fourier series.

For completeness, a continuous periodic function with all Fourier
coefficients zero is zero.  Indeed, the kernel
\[
 K_N(u)=\frac1N\left|\sum_{j=0}^{N-1}e^{2\pi iju}\right|^2
\]
is nonnegative, has integral one over a period, and is a finite Fourier sum.
Away from the integers it is at most $1/(N\sin^2\pi u)$.
Uniform continuity and this bound show that convolution with $K_N$ converges
uniformly to any continuous periodic function.
When all coefficients vanish, each convolution is zero.
Apply this uniqueness statement to $P_x$ minus its displayed Fourier series,
then set $u=0$ to obtain \eqref{eq:theta-reflection}.
Finally
$\theta(x)-1=2e^{-\pi x}\sum_{n\geq1}e^{-\pi(n^2-1)x}$,
and the last sum is bounded by its value at $x=1$.
\end{proof}

\Needspace{5\baselineskip}
\begin{theorem}[Continuation and reflection]
\label{thm:completion}
The function $\mathcal L(s)=\Gamma_{\R}(s)\zeta(s)$, initially defined for
$\operatorname{Re}s>1$, extends meromorphically to $\C$ by
\begin{equation}
 \mathcal L(s)=-\frac1s+\frac1{s-1}
   +\frac12\int_1^\infty(\theta(x)-1)
                   \left(x^{s/2}+x^{(1-s)/2}\right)\frac{dx}{x}.
 \label{eq:theta-mellin}
\end{equation}
It has simple poles at zero and one, with residues $-1$ and $1$.
It satisfies $\mathcal L(s)=\mathcal L(1-s)$ and
$\mathcal L(\bar s)=\overline{\mathcal L(s)}$.
The function
\begin{equation}
 X(s)=\frac12 s(s-1)\mathcal L(s)
 \label{eq:entire-completion}
\end{equation}
is entire, obeys the same symmetries, and satisfies
$X(0)=X(1)=1/2$ and
\begin{equation}
 \log\max_{|s|\leq R}|X(s)|\leq C R\log(R+2)\qquad(R\geq2).
 \label{eq:entire-growth}
\end{equation}
Every zero of $X$ lies in $0\leq\operatorname{Re}s\leq1$.
\end{theorem}

\begin{proof}
For $\operatorname{Re}s>1$, absolute convergence gives
\[
 \frac12\int_0^\infty(\theta(x)-1)x^{s/2}\frac{dx}{x}
 =\sum_{n\geq1}\int_0^\infty e^{-\pi n^2x}x^{s/2}\frac{dx}{x}
 =\Gamma_{\R}(s)\zeta(s).
\]
For the integral from zero to one, use
$\theta(x)-1=x^{-1/2}-1+x^{-1/2}(\theta(1/x)-1)$.
The first two terms integrate to $1/(s-1)-1/s$ after the factor $1/2$.
Substitution $y=1/x$ in the third gives the second term in the integral
of \eqref{eq:theta-mellin}.

The exponential bound in Lemma~\ref{lem:theta} makes that integral entire:
on every compact set, powers of $x$ and any fixed powers of $\log x$ are
bounded by an integrable exponential majorant.
The pole, residue, and symmetry assertions now follow directly.
In particular
\[
 X(s)=\frac12+\frac{s(s-1)}4
       \int_1^\infty(\theta(x)-1)
             \left(x^{s/2}+x^{(1-s)/2}\right)\frac{dx}{x}.
\]
For $|s|\leq R$, the integral is bounded by a constant times
$\int_1^\infty e^{-\pi x}x^{(R+1)/2}\,dx$.
If $r=(R+1)/2$, then
$x^r e^{-\pi x/2}\leq(2r/(\pi e))^r$, by maximizing its logarithm.
The remaining exponential is integrable.  This proves
\eqref{eq:entire-growth}, including the polynomial prefactor.

For $\operatorname{Re}s>1$, the gamma factor and the prime product are
nonzero, so $X$ is nonzero there.  Reflection excludes zeros for
$\operatorname{Re}s<0$.
\end{proof}

The continuation of the original Dirichlet series is now the specific
meromorphic function $\zeta(s)=\pi^{s/2}E(s/2)\mathcal L(s)$.
It has a simple pole at one.  At zero the zero of $E(s/2)$ cancels the
pole of $\mathcal L$ and gives $\zeta(0)=-1/2$.
At each negative even integer, $E(s/2)$ has a simple zero and
$\mathcal L(s)=\mathcal L(1-s)\ne0$.
These are its trivial zeros.  Its other zeros are precisely the zeros of
$X$, with the same multiplicities.  We call them the nontrivial zeros.
No assertion about their real parts inside the closed strip has been used.

The closed strip can be sharpened by a positive trigonometric identity.
The argument uses the pole of $\zeta$ at one and holomorphy at every
other point of the line $\operatorname{Re}s=1$.

\Needspace{5\baselineskip}
\begin{lemma}[The pole and the boundary germs]
\label{lem:boundary-germs}
The continuation
$\zeta(s)=\pi^{s/2}E(s/2)\mathcal L(s)$ is holomorphic on
$\C\setminus\{1\}$ and has a simple pole of residue one at $s=1$.
For every nonzero real $t$, it is holomorphic on neighborhoods of
$1+it$ and $1+2it$.
\end{lemma}
\begin{proof}
Theorem~\ref{thm:completion} makes $\mathcal L$ holomorphic away from
zero and one.  Lemma~\ref{lem:gamma} makes $E$ entire, with
$E(z)=z+O(z^2)$ at zero.
Thus $\pi^{s/2}E(s/2)\mathcal L(s)=-1/2+O(s)$ there, and its
singularity at zero is removable.
The Gaussian integral gives
\[
 \Gamma(1/2)=2\int_0^\infty e^{-u^2}\,du=\sqrt\pi.
\]
Consequently $\pi^{1/2}E(1/2)=1$.
The residue of $\mathcal L$ at one is one, so
\begin{equation}
 \zeta(s)=\frac1{s-1}+H(s)
 \label{eq:boundary-pole}
\end{equation}
on a neighborhood of one, with $H$ holomorphic.
The asserted local holomorphy follows because $t\ne0$ places both
$1+it$ and $1+2it$ away from the only pole.
\end{proof}

\Needspace{5\baselineskip}
\begin{lemma}[A positive prime sum]
\label{lem:boundary-prime}
For every real $\sigma>1$ and real $t$,
\begin{equation}
 \zeta(\sigma)^3
 |\zeta(\sigma+it)|^4|\zeta(\sigma+2it)|\geq1.
 \label{eq:boundary-product}
\end{equation}
\end{lemma}
\begin{proof}
On $\operatorname{Re}s>1$, define the holomorphic function
\[
 B(s)=\sum_p\sum_{r\geq1}\frac{p^{-rs}}r.
\]
Here $p$ runs over positive primes and $p^{-rs}=e^{-rs\log p}$.
For every $\delta>0$, its series converges absolutely and uniformly on
$\operatorname{Re}s\geq1+\delta$, since
\[
 \sum_p\sum_{r\geq1}\frac{p^{-r(1+\delta)}}r
 \leq\frac1{1-2^{-1-\delta}}\sum_{n\geq2}n^{-1-\delta}<\infty.
\]
Lemma~\ref{lem:euler} gives $e^{B(s)}=\zeta(s)$.
Taking absolute values yields $\operatorname{Re}B(s)=\log|\zeta(s)|$,
with the real logarithm of a positive number.
This operation needs no choice of a logarithm at a boundary point.
Absolute convergence permits the three series to be combined.
The identity
\[
 3+4\cos\theta+\cos(2\theta)=2(1+\cos\theta)^2\geq0
 \qquad(\theta\in\R)
\]
therefore gives
\begin{align*}
 &3\log\zeta(\sigma)+4\log|\zeta(\sigma+it)|
                      +\log|\zeta(\sigma+2it)|\\
 &\quad=\sum_p\sum_{r\geq1}\frac{p^{-r\sigma}}r
       \bigl(3+4\cos(rt\log p)+\cos(2rt\log p)\bigr)\geq0.
\end{align*}
The Dirichlet series makes $\zeta(\sigma)>0$.
Exponentiation proves \eqref{eq:boundary-product}.
\end{proof}

The exponents in \eqref{eq:boundary-product} decide the boundary case.
A simple pole contributes order $-3$, while even a simple zero in the
middle factor contributes order $4$.
The third factor is bounded at the limiting point whenever $t\ne0$.

\Needspace{5\baselineskip}
\begin{theorem}[Nonvanishing on the boundary and the strict strip]
\label{thm:strict-strip}
For every real $t\ne0$, one has $\zeta(1+it)\ne0$.
At $s=1$, the function $\zeta$ has a simple pole of residue one.
Every zero $\rho$ of $X$ satisfies
\begin{equation}
 0<\operatorname{Re}\rho<1,
 \qquad
 z_\rho=-i(\rho-1/2),
 \qquad
 |\operatorname{Im}z_\rho|<\frac12.
 \label{eq:strict-strip}
\end{equation}
The zeros of $X$ are exactly the nontrivial zeros of $\zeta$, with
the same multiplicities.
Reflection and complex conjugation preserve these multiplicities.
\end{theorem}
\begin{proof}
Fix $t\in\R\setminus\{0\}$ and suppose $\zeta(1+it)=0$.
By Lemma~\ref{lem:boundary-germs}, this is a zero of a holomorphic
function.  It has some finite order $m\geq1$.
Indeed, a zero of infinite order would make the function vanish on a
disc, contrary to its nonzero Euler product in the same disc's
intersection with $\operatorname{Re}s>1$.
Taylor expansion gives a holomorphic function $a$, with $a(0)\ne0$,
such that
\[
 \zeta(1+it+w)=w^m a(w)
\]
for sufficiently small complex $w$.
For positive real $h$ tending to zero, the pole expansion and local
holomorphy at $1+2it$ give constants $C_0,C_1,C_2>0$ such that
\[
 0<\zeta(1+h)\leq C_0h^{-1},\qquad
 |\zeta(1+h+it)|\leq C_1h^m,\qquad
 |\zeta(1+h+2it)|\leq C_2.
\]
The constants can depend on $t$ and on the chosen neighborhoods.
In particular, boundedness of the third factor allows a zero at $1+2it$.
Lemma~\ref{lem:boundary-prime} now implies
\[
 1\leq\zeta(1+h)^3|\zeta(1+h+it)|^4|\zeta(1+h+2it)|
 \leq C_0^3C_1^4C_2 h^{4m-3}\longrightarrow0.
\]
This contradiction proves nonvanishing for $t\ne0$.
When $t=0$, Lemma~\ref{lem:boundary-germs} gives the pole directly.
Thus the preceding limiting argument never treats the pole as a
holomorphic value of $\zeta$.

For $s=1+it$ with $t\ne0$, the factors $s$, $s-1$, and
$\Gamma_{\R}(s)=\pi^{-s/2}\Gamma(s/2)$ are finite and nonzero.
The nonvanishing just proved therefore gives $X(1+it)\ne0$.
The value $X(1)=1/2$ handles $t=0$.
On the other boundary, the exact reflection identity gives
$X(it)=X(1-it)\ne0$ for every real $t$, including $X(0)=1/2$.
The closed-strip statement in Theorem~\ref{thm:completion} now yields
$0<\operatorname{Re}\rho<1$ for every zero $\rho$ of $X$.

On this open strip, $s(s-1)\Gamma_{\R}(s)/2$ is holomorphic and
nonzero.  Multiplication by it preserves each zero order, so the zeros
of $X$ and $\zeta$ there have the same multiplicities.
At $s=0$, one has $\zeta(0)=-1/2$.
For $\operatorname{Re}s<0$, reflection gives
$\mathcal L(s)=\mathcal L(1-s)\ne0$.
The reciprocal gamma factor has simple zeros exactly at
$s=0,-2,-4,\ldots$.
It follows that the zeros of $\zeta$ in $\operatorname{Re}s<0$ are
exactly the simple zeros at $s=-2,-4,\ldots$.
These are the trivial zeros.
On $\operatorname{Re}s=0$ with $s\ne0$, both $E(s/2)$ and
$\mathcal L(s)$ are nonzero.  The Euler product excludes zeros for
$\operatorname{Re}s>1$.
This proves the stated classification of all other zeros.

If $\rho=\beta+i\gamma$ with $\beta,\gamma\in\R$, then
\[
 z_\rho=\gamma+i(1/2-\beta).
\]
Thus $0<\beta<1$ is exactly the strict bound in
\eqref{eq:strict-strip}.
The affine map $s\mapsto-i(s-1/2)$ has nonzero derivative.
Consequently a zero of order $m$ at $s=\rho$ is a zero of order $m$
of $z\mapsto X(1/2+iz)$ at $z=z_\rho$.
Finally the identities $X(s)=X(1-s)$ and
$X(\bar s)=\overline{X(s)}$ preserve orders in the local power series.
In centered variables their images are
\[
 z_{1-\rho}=-z_\rho,\qquad
 z_{\bar\rho}=-\overline{z_\rho},\qquad
 z_{1-\bar\rho}=\overline{z_\rho}.
\]
All three parameters therefore occur with the same multiplicity as
$z_\rho$, with coincident parameters counted only according to their
actual zero order.
\end{proof}

Lemma~\ref{lem:test-bounds} supplies exponential decay for every test in $\A$.
The same centered coordinate specifies where the Weil transform is
evaluated.  Suppose $a>1/2$ and $g:\R\to\C$ is continuous with
$|g(u)|\leq C e^{-a|u|}$.
Define
\[
 \widehat g(z)=\int_{\R}g(u)e^{izu}\,du,
 \qquad
 F_g(s)=\widehat g(-i(s-1/2))
       =\int_{\R}g(u)e^{(s-1/2)u}\,du.
\]
For each $0\leq b<a$ and integer $k\geq0$, the absolute value of the
$k$th derivative integrand is bounded on $|\operatorname{Im}z|\leq b$ by
\[
 C|u|^k e^{-(a-b)|u|}.
\]
This is integrable on $\R$.
Termwise differentiation therefore makes $\widehat g$ holomorphic on
$|\operatorname{Im}z|<a$.
It makes $F_g$ holomorphic on
$1/2-a<\operatorname{Re}s<1/2+a$.
Every nontrivial zero lies in this domain, and
$F_g(\rho)=\widehat g(z_\rho)$ uses the same multiplicity under the
invertible affine change of variable.

The strict strip permits zeros with $\operatorname{Re}\rho\ne1/2$.
The assertion that all $z_\rho$ are real is the Riemann hypothesis.
It remains an additional assertion beyond boundary nonvanishing.

\Needspace{5\baselineskip}
\begin{corollary}[Zero count and absolute test sums]
\label{cor:zero-count}
The number of nontrivial zeros with $|\rho|\leq R$, counted with multiplicity,
is $O(R\log(R+2))$.
For every $g\in\A$, the series $\sum_\rho|F_g(\rho)|$ converges.
\end{corollary}

\begin{proof}
Apply Lemma~\ref{lem:jensen} to $X$, using $X(0)=1/2$.
Since $0\leq\operatorname{Re}\rho\leq1$, Lemma~\ref{lem:test-bounds}
gives $|F_g(\rho)|\leq C_N(1+|\operatorname{Im}\rho|)^{-N}$ for every $N$.
The zeros in $2^j\leq|\operatorname{Im}\rho|<2^{j+1}$ contribute at most
$C(j+1)2^{j(1-N)}$.  Taking $N=3$ gives a convergent series.
The bounded region has only finitely many zeros.
\end{proof}

\section{Heights at which the horizontal integrals vanish}
\label{sec:heights}

Write $H(s)=\mathcal L'(s)/\mathcal L(s)$.
It has residue equal to the multiplicity at each nontrivial zero.
Its residue at each of zero and one is $-1$.
These signs follow by differentiating the logarithm of a local factor
$(s-s_0)^m$ or $(s-s_0)^{-1}$.
The identity
\begin{equation}
 H(s)=\frac{X'(s)}{X(s)}-\frac1s-\frac1{s-1},
 \qquad H(1-s)=-H(s)
 \label{eq:log-reflection}
\end{equation}
will control both the poles and the left vertical line.

\Needspace{5\baselineskip}
\begin{lemma}[Admissible rectangle heights]
There are heights\label{lem:heights} $T_m\in[m,m+1]$, for all sufficiently large integers $m$,
such that no zero lies on either horizontal segment
$-1/2\leq\operatorname{Re}s\leq3/2$, $\operatorname{Im}s=\pm T_m$,
and
\begin{equation}
 \sup_{-1/2\leq\sigma\leq3/2}|H(\sigma\pm iT_m)|
       =O(T_m^2\log^2(T_m+2)).
 \label{eq:horizontal-H}
\end{equation}
For $g\in\A_a$ and $1<c<\min(a+1/2,3/2)$, both horizontal integrals
of $H(s)F_g(s)$ between the lines $1-c$ and $c$ tend to zero at these heights.
\end{lemma}

\begin{proof}
Let $N_m$ count all zeros in $|s|\leq m+8$, with multiplicity, and set
$\delta_m=1/(8(N_m+1))$.
Remove from $[m,m+1]$ the intervals of radius $\delta_m$ about the absolute
values of their imaginary parts.
The removed length is at most $2\delta_mN_m<1/4$, so a remaining $T_m$ exists.
Zeros outside this disc have $|\operatorname{Im}\rho|>m+1+\delta_m$
for large $m$, because $0\leq\operatorname{Re}\rho\leq1$.
Thus both heights have distance at least $\delta_m$ from every zero ordinate.
Corollary~\ref{cor:zero-count} gives
$\delta_m^{-1}=O(m\log(m+2))$.

Fix $T=T_m$ and center a disc at $s_0=2+iT$.
By \eqref{eq:entire-growth},
$|X(s_0+z)|\leq\exp(C(T+8)\log(T+8))$ for $|z|\leq6$.
At its center, the prime-product lower bound and
\eqref{eq:reciprocal-growth} give
\[
 |X(2+iT)|
 =\frac{|(2+iT)(1+iT)|}{2\pi}
       |\Gamma(1+iT/2)|\,|\zeta(2+iT)|
 \geq\exp(-C(T+8)\log(T+8)).
\]
Choose $R\in[5,6]$ with no zero on the circle.
The number of zeros in it is $O(T\log(T+2))$ by the global count.
Every point of the desired horizontal segment has $|s-s_0|\leq5/2$,
and its distance from every zero is at least $\delta_m$.
Lemma~\ref{lem:local-log} now gives
\[
 |X'(s)/X(s)|
 \leq C T\log(T+2)(1+\delta_m^{-1})
 =O(T^2\log^2(T+2)).
\]
The two rational terms in \eqref{eq:log-reflection} are bounded at these
heights.  Conjugation gives the same estimate on the bottom segment.

The chosen vertical strip lies strictly inside
$1/2-a<\operatorname{Re}s<1/2+a$.
Its test transform is $O(T^{-4})$ uniformly by
Lemma~\ref{lem:test-bounds}.
The horizontal lengths are bounded, so their integrals are
$O(T^{-2}\log^2(T+2))$, which tends to zero.
\end{proof}

The proof uses a polynomial bound, which suffices for every test in $\A$.
It does not require a sharper local zero count, a product over all zeros,
or an asymptotic formula for the number of zeros.

\Needspace{7\baselineskip}
\section{The normalized explicit formula}
\label{sec:formula}

\begin{theorem}[The full arithmetic functional]
\label{thm:explicit}
For every $g\in\A$, equation~\eqref{eq:target} holds, with all sums and
integrals absolutely convergent.
Its right side defines a linear functional $L:\A\to\C$ satisfying
\begin{equation}
 L(g^*)=\overline{L(g)},\qquad
 L(g)=\sum_\rho\widehat g\bigl(-i(\rho-1/2)\bigr).
 \label{eq:functional}
\end{equation}
\end{theorem}

\begin{proof}
Choose $g\in\A_a$ and $1<c<\min(a+1/2,3/2)$.
The prime terms converge because Lemma~\ref{lem:test-bounds} gives
\[
 \sum_{n\geq2}\frac{\Lambda_{\mathrm{vM}}(n)}{\sqrt n}
                    (|g(\log n)|+|g(-\log n)|)
 \leq C\sum_{n\geq2}(\log n)n^{-a-1/2}<\infty.
\]
The gamma integral converges by \eqref{eq:digamma} and
\eqref{eq:strip-decay}.  The polar values are inside the transform strip,
and the zero sum converges by Corollary~\ref{cor:zero-count}.

Integrate $H(s)F_g(s)$ around the rectangle with vertical sides $1-c,c$
and heights $\pm T_m$, counterclockwise.  Its right side is traversed upward.
There are no poles on the vertical sides, since the zeros lie in the
closed critical strip.
The residue formula gives the zero sum inside the rectangle minus
$F_g(0)+F_g(1)$.
Lemma~\ref{lem:heights} eliminates the horizontal integrals.
On $\operatorname{Re}s=c$, equations~\eqref{eq:euler} and
\eqref{eq:digamma} give $H(c+it)=O(\log(2+|t|))$.
Reflection gives the same bound on $1-c$, so both vertical integrals
converge absolutely.
In the left integral, substitute $s=1-w$ and use $H(1-w)=-H(w)$.
Both vertical-line integrals in the resulting expression are oriented upward.
The outcome is
\begin{equation}
 \sum_\rho F_g(\rho)-F_g(0)-F_g(1)
 =\frac1{2\pi i}\int_{\operatorname{Re}s=c}
           H(s)\bigl(F_g(s)+F_g(1-s)\bigr)\,ds.
 \label{eq:spine-contour}
\end{equation}

On this line, split $H=H_\Gamma+\zeta'/\zeta$.
The series for $\zeta'/\zeta$ can pass through the integral because
$\sum\Lambda_{\mathrm{vM}}(n)n^{-c}<\infty$ and both test factors are integrable
on the vertical line.
For the first factor, Fourier inversion applied to
$u\mapsto e^{(c-1/2)u}g(u)$ gives
\begin{equation}
 \frac1{2\pi i}\int_{\operatorname{Re}s=c}n^{-s}F_g(s)\,ds
 =n^{-c}e^{(c-1/2)\log n}g(\log n)
 =n^{-1/2}g(\log n).
 \label{eq:prime-inversion}
\end{equation}
This weighted function has exponential decay with all derivatives, since
$c-1/2<a$.  The proof of Lemma~\ref{lem:gaussian} therefore applies to it.
For $F_g(1-s)$, use $u\mapsto e^{(c-1/2)u}g(-u)$ instead.
It gives $n^{-1/2}g(-\log n)$.
The minus sign in \eqref{eq:euler} proves exactly the prime-power term in
\eqref{eq:target}.

It remains to evaluate the gamma contribution.
The function $H_\Gamma$ is holomorphic for $\operatorname{Re}s>0$.
Shift each gamma integral from $c$ to $1/2$.
There is no pole between these lines, and the logarithmic bound
\eqref{eq:digamma} times the rapid test decay eliminates the horizontal sides.
Set $s=1/2+it$.  The two factors are $\widehat g(t)$ and $\widehat g(-t)$.
Changing $t$ to $-t$ in the second term gives
\begin{align*}
 &\frac1{2\pi}\int_{\R}\widehat g(t)
         \bigl(H_\Gamma(1/2+it)+H_\Gamma(1/2-it)\bigr)\,dt\\
 &\hspace{1cm}=\frac1{2\pi}\int_{\R}\widehat g(t)
     \left[\operatorname{Re}\psi(1/4+it/2)-\log\pi\right]dt.
\end{align*}
Here conjugation of $\Gamma$ follows from its defining real integral and
analytic continuation.  Together with \eqref{eq:spine-contour}, this proves
\eqref{eq:target}.

Finally the gamma bracket is real and even.
For real $t$, \eqref{eq:star-transform} gives
$\widehat{g^*}(t)=\overline{\widehat g(t)}$.
The involution interchanges the two polar values and the two values at
$\pm\log n$, followed by conjugation.
Thus all three local terms obey the Hermitian identity in
\eqref{eq:functional}.
\end{proof}

The proof explains the normalization directly.
The poles of $\mathcal L$ each contribute residue $-1$ to its logarithmic
derivative, giving the two positive polar values after rearrangement.
The logarithmic derivative of each Euler factor gives the negative prime
term.  The reflection combines the two gamma factors at $1/2\pm it$.
These operations preserve the original complex test function throughout.

\section{Gaussian tests and their boundary terms}
\label{sec:polynomial-gaussians}

For $g_A$ from Lemma~\ref{lem:gaussian}, substitution in the formula gives
\begin{equation}
\begin{split}
 \sum_\rho e^{A(\rho-1/2)^2}
 ={}&2e^{A/4}
   +\frac1{2\pi}\int_{\R}e^{-At^2}
      \left[\operatorname{Re}\psi(1/4+it/2)-\log\pi\right]dt\\
 &-\frac1{\sqrt{\pi A}}\sum_{n\geq2}
       \frac{\Lambda_{\mathrm{vM}}(n)}{\sqrt n}
                       e^{-(\log n)^2/(4A)}.
\end{split}
\label{eq:gaussian-explicit}
\end{equation}
If every nontrivial zero has real part $1/2$, each term in this Gaussian
zero sum is positive.  The identity itself is unconditional.

\Needspace{5\baselineskip}
\begin{proposition}[Polynomial--Gaussian compatibility]
\label{prop:polynomial-gaussian}
Let $A>0$, $\tau\in\R$, and $P\in\C[w]$.  The function
\begin{equation}
 f(u)=e^{-i\tau u}P(i\partial_u)g_A(u)
 \label{eq:gaussian-inverse}
\end{equation}
belongs to $\A_a$ for every $a>0$, and
\begin{equation}
 \widehat f(z)=P(z-\tau)e^{-A(z-\tau)^2}.
 \label{eq:polynomial-transform}
\end{equation}
For every fixed $b>0$ and $z=t+iy$ with $|y|\leq b$,
\begin{equation}
 |\widehat f(z)|\leq C_{P,A,b,\tau}(1+|t|)^{\deg P}
                             e^{-A(t-\tau)^2+A b^2}.
 \label{eq:gaussian-strip}
\end{equation}
Consequently every real integration-by-parts boundary value and every
horizontal contour term used in Theorem~\ref{thm:explicit} vanishes for
$f$ and $f^*\star f$.
\end{proposition}

\begin{proof}
Each derivative in \eqref{eq:gaussian-inverse} is a polynomial in $u$ times
$e^{-u^2/(4A)}$, multiplied by $e^{-i\tau u}$.
This description holds after any further derivative.
Multiplication by $e^{a|u|}$ still leaves an integrable function tending
to zero at both infinities, for every $a>0$.
Integration by parts gives
$\widehat{i\partial_u h}(z)=z\widehat h(z)$, and modulation gives
$\widehat{e^{-i\tau u}h}(z)=\widehat h(z-\tau)$.
Together these prove \eqref{eq:polynomial-transform}.
The exact Gaussian modulus is
$|e^{-A(t-\tau+iy)^2}|=e^{-A(t-\tau)^2+Ay^2}$,
which proves \eqref{eq:gaussian-strip}.
This bound times \eqref{eq:horizontal-H} tends to zero at both heights.
The convolution square lies in every $\A_a$ by
Proposition~\ref{prop:algebra}; its transform is the product in
\eqref{eq:star-transform}, with Gaussian decay of twice the exponent.
Thus it has the same boundary properties.
\end{proof}

\Needspace{5\baselineskip}
\begin{example}[The first polynomial and a missing conjugation]
For $P(w)=w$, the function and its transform are
\[
 \begin{aligned}
 f(u)&=-iu\,e^{-i\tau u}g_A(u)/(2A),\\
 \widehat f(z)&=(z-\tau)e^{-A(z-\tau)^2}.
 \end{aligned}
\]
Its convolution square has transform
\[
 \widehat{f^*\star f}(z)=(z-\tau)^2e^{-2A(z-\tau)^2}.
\]
At $z=\tau+ib$, this value is $-b^2e^{2Ab^2}$.
The correct analytic expression is
$\overline{\widehat f(\bar z)}\widehat f(z)$.
Replacing it by $|\widehat f(z)|^2$ would erase the negative sign and would
not give a holomorphic function of $z$.
\end{example}

\section{A separate positivity argument}
\label{sec:positivity}

The explicit formula compares linear functionals.
Positivity concerns its value on every convolution square and requires a
different argument.  Define the Riemann hypothesis to mean that every
nontrivial zero $\rho$ satisfies $\operatorname{Re}\rho=1/2$.

\Needspace{5\baselineskip}
\begin{theorem}[Positivity on the exact test algebra]
\label{thm:positivity}
The Riemann hypothesis is equivalent to
\begin{equation}
 L(f^*\star f)\geq0\qquad\text{for every }f\in\A.
 \label{eq:positivity}
\end{equation}
Under the Riemann hypothesis, writing $\rho=1/2+i\gamma_\rho$ gives
$L(f^*\star f)=\sum_\rho|\widehat f(\gamma_\rho)|^2$.
If there is a zero off the critical line, a function of the form
\eqref{eq:gaussian-inverse} has strictly negative convolution-square value.
\end{theorem}

\begin{proof}
Put $z_\rho=-i(\rho-1/2)$.
Theorem~\ref{thm:completion} gives $|\operatorname{Im}z_\rho|\leq1/2$.
Its reflection and conjugation symmetries show that $z_\rho$ and
$\bar z_\rho$ occur with equal multiplicity.
Under the Riemann hypothesis all these parameters are real.
Then Theorem~\ref{thm:explicit} and \eqref{eq:star-transform} give the
asserted nonnegative sum.

For the converse, suppose $z_0=\tau+ib$ occurs, where $b\ne0$.
Its conjugate has the same multiplicity $m$.
Center the parameter by $w=z-\tau$.
There are finitely many distinct centered zero parameters with
$|\operatorname{Re}w|\leq2$.
This follows from the global zero count and the bound on their imaginary
parts.  Form the monic polynomial $Q$ whose roots are these distinct points
except $ib$ and $-ib$.  If there are none, set $Q=1$.
This root set is invariant under conjugation, so $Q$ has real coefficients
and $Q(ib)Q(-ib)=|Q(ib)|^2\ne0$.
Set
\[
 P(w)=wQ(w)Q(-w),\qquad
 \widehat f_A(z)=P(z-\tau)e^{-A(z-\tau)^2},\qquad A\geq1.
\]
The polynomial $P$ is real and odd.
Proposition~\ref{prop:polynomial-gaussian} constructs $f_A\in\A$ exactly.
Since $\tau$ and the coefficients are real,
$\overline{\widehat f_A(\bar z)}=\widehat f_A(z)$.
The two selected zero parameters therefore contribute
\begin{equation}
 -2m b^2|Q(ib)|^4e^{2Ab^2}
 \label{eq:negative-pair}
\end{equation}
to $L(f_A^*\star f_A)$.
All other centered zero parameters with $|\operatorname{Re}w|\leq2$
give zero because $Q$ vanishes there.

For a remaining parameter $w=x+iy$, one has $|x|>2$, $|y|\leq1/2$,
and $D=x^2-y^2\geq15/4$.
For $A\geq1$, $2AD\geq15A/4+(x^2-1/4)$.
Thus the absolute value of the remaining sum is bounded by
\begin{equation}
 e^{-15A/4}
 \sum_{|\operatorname{Re}w|>2}|P(w)|^2e^{-(x^2-1/4)}.
 \label{eq:separator-tail}
\end{equation}
The final series is finite by the polynomial growth of $P$ and the zero
count.  It is independent of $A$.
Expression~\eqref{eq:negative-pair} tends to minus infinity while
\eqref{eq:separator-tail} tends to zero.
The Hermitian identity in \eqref{eq:functional} makes the full value real,
so it is strictly negative for sufficiently large $A$.
This contradicts \eqref{eq:positivity} and proves the converse.
\end{proof}

The conclusion is a criterion, not a proof of the Riemann hypothesis.
The transform estimates, the continuation, and the contour formula were
proved before invoking any sign condition.
Their exact comparison identifies the remaining arithmetic question:
whether the full local expression is nonnegative on every convolution square.
The negative separator describes what any off-line zero would force.


\section{The form on the nonunital test algebra}
\label{sec:native-form}

For the representation theory, write
$\cA_\zeta=\A$ and $L_\zeta=L$.
These are identical sets, algebra operations, and functional values.
This identification takes no completion.  Membership in $\A$ means that there is
$a>1/2$ with every $p_{a,k}$ finite.  Setting $\epsilon=a-1/2$ says precisely
that $e^{(1/2+\epsilon)|u|}g^{(k)}(u)$ is integrable for every $k$.
The involution is the reflection and conjugation in
\eqref{eq:operations}; its square transform is the analytic product
\eqref{eq:star-transform}.

The operations on this algebra are
\begin{equation}
 (f\star g)(u)=\int_{\R}f(v)g(u-v)\,dv,
 \qquad f^*(u)=\overline{f(-u)}.
 \label{eq:conv}
\end{equation}
On each common holomorphic strip their transforms obey
\begin{equation}
 \widehat{f\star g}(z)=\widehat f(z)\widehat g(z),
 \qquad\widehat{f^*}(z)=\overline{\widehat f(\bar z)}.
 \label{eq:transform}
\end{equation}
These are the operations and identities already proved in
\eqref{eq:operations} and Proposition~\ref{prop:algebra}.
For real $u$, the same Fourier convention and Lemma~\ref{lem:gaussian} give
\begin{equation}
 \widehat f(z)=\int_{\R}f(u)e^{izu}\,du,
 \qquad f(u)=\frac1{2\pi}\int_{\R}\widehat f(t)e^{-itu}\,dt.
 \label{eq:fourier}
\end{equation}
In particular, the product identity on the real Fourier axis is
\begin{equation}
 \widehat{f^*\star f}(t)=|\widehat f(t)|^2\qquad(t\in\R).
 \label{eq:square-transform}
\end{equation}
The derivative of the gamma normalization in \eqref{eq:gamma-real} reads
\begin{equation}
 2\operatorname{Re}\frac{\Gamma_{\R}'}{\Gamma_{\R}}(1/2+it)
 =\operatorname{Re}\psi(1/4+it/2)-\log\pi.
 \label{eq:gamma-normalization}
\end{equation}

The multiplicative variable $x=e^u$ gives the identity
\begin{equation}
 \widehat g\bigl(-i(s-1/2)\bigr)
 =\int_0^\infty g(\log x)x^{s-1/2}\,\frac{dx}{x},
 \qquad 1/2-a<\operatorname{Re}s<1/2+a
 \label{eq:mellin}
\end{equation}
for $g\in\A_a$.
The centered Mellin convention gives
$F_g(0)=\widehat g(i/2)$ and $F_g(1)=\widehat g(-i/2)$.
Thus the functional constructed in Theorem~\ref{thm:explicit} is exactly
\begin{equation}
\begin{split}
 L_\zeta(g)={}&\widehat g(i/2)+\widehat g(-i/2)\\
 &+\frac1{2\pi}\int_{\R}\widehat g(t)
   \left[\operatorname{Re}\psi(1/4+it/2)-\log\pi\right]dt\\
 &-\sum_{n\geq2}\frac{\Lambda_{\mathrm{vM}}(n)}{\sqrt n}
                 \bigl(g(\log n)+g(-\log n)\bigr).
\end{split}
\label{eq:weil-functional}
\end{equation}
All sums and integrals converge absolutely on this algebra.
Theorem~\ref{thm:explicit} also gives the exact zero expression
\begin{equation}
 L_\zeta(g)=\sum_\rho\widehat g\bigl(-i(\rho-1/2)\bigr).
 \label{eq:explicit}
\end{equation}
A sesquilinear form is linear in its first argument and conjugate-linear in
its second.  It is Hermitian when reversing the arguments conjugates the
value, and positive when its diagonal values are nonnegative real numbers.
Define the sesquilinear form by
\[
 B_\zeta(f,g)=L_\zeta(g^*\star f).
\]
The identity
\begin{equation}
 L_\zeta(g^*)=\overline{L_\zeta(g)}
 \label{eq:hermitian}
\end{equation}
from Theorem~\ref{thm:explicit} makes $B_\zeta$ Hermitian.
Theorem~\ref{thm:positivity} states
\begin{equation}
 \mathrm{RH}\quad\Longleftrightarrow\quad
 L_\zeta(f^*\star f)\geq0\quad\hbox{for every }f\in\cA_\zeta.
 \label{eq:weil-criterion}
\end{equation}
Under RH, the same theorem gives
\begin{equation}
 L_\zeta(f^*\star f)=\sum_\rho|\widehat f(\gamma_\rho)|^2,
 \qquad \rho=1/2+i\gamma_\rho.
 \label{eq:zero-norm}
\end{equation}
The right side explains positivity conditional on RH.  Both the reality
of the parameters $\gamma_\rho$ and their occurrence in this expression
use the stated zero placement.  They do not construct positivity from the
arithmetic local terms independently.

Theorem~\ref{thm:strict-strip} proves
$0<\operatorname{Re}\rho<1$ and $|\operatorname{Im}z_\rho|<1/2$,
with the same zero multiplicities in both coordinates.
The closed bounds used in the contour and separator follow from these
strict bounds.  Nonreal $z_\rho$ are still permitted, so boundary
nonvanishing supplies no independent positivity of $B_\zeta$.

\section{Positive functionals and their representations}

A $*$-homomorphism preserves addition, complex scalar multiplication,
products, and this involution; no preservation of a unit is imposed here.
For a positive form, expansion of the diagonal at $x+\lambda y$ gives
Hermitian symmetry and
\[
 |\langle x,y\rangle|^2\leq\langle x,x\rangle\langle y,y\rangle.
\]
Indeed, when $\langle y,y\rangle>0$, minimize that quadratic expression in
$\lambda$; if it is zero, arbitrarily large multiples force the cross term
to vanish.  A positive definite form gives a norm
$\|x\|=\sqrt{\langle x,x\rangle}$.  Its space is a pre-Hilbert space.
Completing Cauchy sequences modulo sequences tending to zero gives a
Hilbert space; continuity extends the inner product to that completion.
An isometry preserves this inner product, and a surjective isometry is
unitary.  An operator is bounded if $\|Sx\|\leq C\|x\|$ on its domain;
the least such constant is its operator norm.

Orthogonality means vanishing of the inner product.  For a linear subspace
$M$ of a Hilbert space $H$, write
\[
 M^\perp=\{v\in H:\langle v,m\rangle=0\text{ for every }m\in M\}.
\]
Closed subspaces have the following decomposition, which will distinguish
a dense operator range from a proper one.

\begin{lemma}[Projection onto a closed subspace]
\label{lem:weil-orthogonal-decomposition}
If $M\subset H$ is a closed linear subspace, every $v\in H$ has a unique
expression $v=m+n$ with $m\in M$ and $n\in M^\perp$.
In particular, $M^\perp=\{0\}$ implies $M=H$.
\end{lemma}
\begin{proof}
Put $d^2=\inf_{m\in M}\|v-m\|^2$ and choose $m_j\in M$ with
$\|v-m_j\|^2\to d^2$.  Expansion of inner products gives the
parallelogram identity and hence
\[
 \|m_j-m_k\|^2
 =2\|v-m_j\|^2+2\|v-m_k\|^2
      -4\left\|v-\frac{m_j+m_k}{2}\right\|^2
 \leq2\|v-m_j\|^2+2\|v-m_k\|^2-4d^2.
\]
Thus $(m_j)$ is Cauchy.  Completeness and closedness give a limit
$m\in M$ attaining the infimum.
For $y\in M$, the real function $t\mapsto\|v-m-ty\|^2$ is minimized
at zero, so its linear coefficient gives
$\operatorname{Re}\langle v-m,y\rangle=0$.
Replacing $y$ by $iy$ also makes the imaginary part zero.
Therefore $n=v-m$ lies in $M^\perp$.
The intersection $M\cap M^\perp$ is zero, since an element of that
intersection has squared norm zero.  This proves uniqueness and the
last assertion.
\end{proof}

An operator on a Hilbert space $H$ means a linear map $S:D\to H$ on a
specified dense linear subspace.  Its graph is $\{(x,Sx):x\in D\}$.
It is closed when its graph is closed in $H\oplus H$, and closable when
the graph closure is again a graph.  Equivalently, $x_n\to0$ and
$Sx_n\to y$ must imply $y=0$; this equivalence follows by subtracting two
possible limiting graph values at the same vector.
The graph norm is $(\|x\|^2+\|Sx\|^2)^{1/2}$.  A core for a closable
operator is a subspace of its domain dense in this norm.
The adjoint has domain
\[
 \mathcal D(S^*)=\{y\in H:\text{there is }z\in H\text{ with }
           \langle Sx,y\rangle=\langle x,z\rangle\text{ for all }x\in D\},
 \qquad S^*y=z.
\]
Density makes $z$ unique.  Passing to limits in this identity shows that
the adjoint is closed.  If this adjoint domain is dense, the preceding
sequence test proves closability: a proposed $y$ pairs to zero with every
adjoint-domain vector.  A symmetric operator satisfies $S\subseteq S^*$;
a self-adjoint operator satisfies equality, including its domain.

A representation on $D$ is an algebra homomorphism into linear maps
$D\to D$.  It is a $*$-representation when
$\langle\pi(a)x,y\rangle=\langle x,\pi(a^*)y\rangle$ for $x,y\in D$.
Thus the algebraic domain must be invariant before the adjoint identity
can be used in repeated products.  We now construct such a representation
from a positive functional; this is the nonunital GNS construction.

Let $\mathfrak A$ be a complex $*$-algebra, with no norm or unit assumed,
and let $W:\mathfrak A\to\C$ be positive:
\[
 W(a^*a)\geq0\qquad(a\in\mathfrak A).
\]

\Needspace{5\baselineskip}
\begin{theorem}[Nonunital algebraic GNS]
\label{thm:gns}
Put
\[
 \langle a,b\rangle_W=W(b^*a),
 \qquad
 \cN_W=\{a:W(a^*a)=0\}.
\]
Then $\cN_W$ is a left ideal, meaning a linear subspace preserved by
left multiplication by every algebra element, and
\begin{equation}
 D_W=\mathfrak A/\cN_W,
 \qquad
 \pi_W(a)[b]=[ab],
 \label{eq:gns-action}
\end{equation}
is a $*$-representation on a pre-Hilbert space.  Regard $D_W$ as a
common invariant dense domain in its Hilbert completion $\mathcal H_W$.
For every $a\in\mathfrak A$,
\begin{equation}
 \pi_W(a^*)\subseteq\pi_W(a)^*,
 \label{eq:adjoint-inclusion}
\end{equation}
so $\pi_W(a)$ is closable.  It extends boundedly to $\mathcal H_W$ if
and only if there is a constant $C_a$ such that
\begin{equation}
 W(b^*a^*ab)\leq C_a^2W(b^*b)
 \qquad(b\in\mathfrak A).
 \label{eq:multiplier-bound}
\end{equation}
The least such $C_a$ is the norm of the extension.
\end{theorem}

\begin{proof}
Positivity and polarization make $\langle\ ,\ \rangle_W$ Hermitian and
give Cauchy--Schwarz.  If $b\in\cN_W$, then
\[
 W((ab)^*ab)=\langle a^*ab,b\rangle_W=0,
\]
so $\cN_W$ is a left ideal and \eqref{eq:gns-action} is well defined.
For classes in $D_W$,
\[
 \langle\pi_W(a)[b],[c]\rangle_W
 =W(c^*ab)
 =\langle[b],\pi_W(a^*)[c]\rangle_W.
\]
This proves \eqref{eq:adjoint-inclusion}.  The adjoint domain contains
the dense subspace $D_W$, so $\pi_W(a)$ is closable.  Finally,
\[
 \|\pi_W(a)[b]\|^2=W(b^*a^*ab),
\]
so the bounded-operator criterion is exactly
\eqref{eq:multiplier-bound}.
\end{proof}

For $a=a^*$, the operator $\pi_W(a)$ is symmetric and closable.
The closure need not be self-adjoint or admit a self-adjoint extension.
The obstruction can be checked directly.  For a closed symmetric operator
$S$, expansion gives
$\|(S\pm i)x\|^2=\|Sx\|^2+\|x\|^2$.
Consequently both ranges are closed, and
\[
 V:(S+i)x\longmapsto(S-i)x
\]
is an isometry from $\operatorname{ran}(S+i)$ onto
$\operatorname{ran}(S-i)$.  Their orthogonal complements are respectively
$\ker(S^*-i)$ and $\ker(S^*+i)$, by the definition of the adjoint.
If a self-adjoint $T$ extends $S$, each $T\pm i$ has closed dense range:
its orthogonal complement is $\ker(T\mp i)=0$ by the norm identity.
Lemma~\ref{lem:weil-orthogonal-decomposition} therefore makes each range
all of $H$.
Thus $(T-i)(T+i)^{-1}$ is a unitary extension of $V$.
Its restrictions to the two range complements must be unitary too, so
these complements must have the same dimension.

For example,
\[
 P=-i\frac{d}{dx},\qquad \mathcal D(P)=C_c^\infty(0,\infty)
 \subset L^2(0,\infty),
\]
Here $C_c^\infty(0,\infty)$ consists of smooth functions with compact
support inside the half-line.  It is dense in $L^2$: truncate the support
and size of an $L^2$ function, approximate by step functions, and replace interval endpoints
by smooth ramps on sets of arbitrarily small length.  The operator is
symmetric by integration by parts, and hence closable by the dense
adjoint-domain criterion.  The adjoint equations, tested against compactly
supported smooth functions, say weakly that $u'=-u$ and $u'=u$.
Multiplication by the integrating factors $e^x$ and $e^{-x}$ gives a weak
zero derivative.  A locally integrable function with weak zero derivative
is constant: every compact test function of integral zero is the derivative
of another compact test function, so the resulting distribution depends
only on the integral of its test.  The solutions are therefore $ce^{-x}$
and $ce^x$, respectively.  Only the first belongs to $L^2(0,\infty)$.
Passing to the closure does not change the adjoint, since its defining
identity passes through graph limits.  Its two range complements thus
have dimensions one and zero, which excludes a self-adjoint extension.

This operator is itself a GNS multiplier.  Put $D=\mathcal D(P)$ and use
on $L^2(0,\infty)$ the inner product linear in its first argument.
For $f,g\in D$, define the rank-one operator
\[
 R_{f,g}x=\langle x,g\rangle f\qquad(x\in D).
\]
Let $\mathfrak A_D$ be the unital algebra of operators $D\to D$
generated by $P$ and all $R_{f,g}$.  Every generator preserves $D$;
integration by parts and the inner product identity give adjoints on $D$
by $I^*=I$, $P^*=P$, and $R_{f,g}^*=R_{g,f}$.  For a finite linear combination of
words, conjugate the coefficients and reverse each word with these
adjoints.  The resulting operator $B^*$ satisfies
$\langle Bx,y\rangle=\langle x,B^*y\rangle$ for $x,y\in D$.
If two expressions define the same operator $B$, this identity and density
of $D$ make their adjoint operators equal.  Thus the involution is
well-defined on $\mathfrak A_D$, and is the restriction of the Hilbert
adjoint to $D$.  It makes $\mathfrak A_D$ a $*$-algebra with invariant
domain $D$.

Choose $\Omega\in D$ with $\|\Omega\|=1$ and put
$W(B)=\langle B\Omega,\Omega\rangle$ for $B\in\mathfrak A_D$.
Then
\[
 W(B^*B)=\|B\Omega\|^2,\qquad
 \mathcal N_W=\{B:B\Omega=0\}.
\]
The map $U_0:[B]\mapsto B\Omega$ from its GNS quotient to $D$ is an
isometry, since
$W(C^*B)=\langle B\Omega,C\Omega\rangle$.
It is onto: $R_{f,\Omega}\Omega=f$ for every $f\in D$.
Its Hilbert completion is therefore all of $L^2(0,\infty)$, and
\[
 U_0\pi_W(C)[B]=CB\Omega=C U_0[B]
 \qquad(B,C\in\mathfrak A_D).
\]
In particular, the symmetric algebra element $P$ acts in this GNS
representation as exactly the operator $P$ with domain $D$.
The domain is invariant under the entire algebra and is a core for
$\overline P$ by the definition of graph closure.  The preceding
range-complement calculation therefore gives a symmetric GNS multiplier
whose closure has no self-adjoint extension.

Products of closures require a separate common-domain construction.
We give that construction in Section~\ref{sec:graph-domain}.

\Needspace{5\baselineskip}
\begin{example}[An unbounded positive GNS multiplier]
Let $\mathfrak A=\C[x]$, let $x^*=x$, and set
\[
 W(p)=\frac1{\sqrt\pi}\int_{\R}p(t)e^{-t^2}\,dt.
\]
Then $W$ is positive, but multiplication by $x$ is unbounded, since
\[
 \frac{\|\pi_W(x)[x^n]\|^2}{\|[x^n]\|^2}
 =\frac{\int_{\R}t^{2n+2}e^{-t^2}\,dt}
        {\int_{\R}t^{2n}e^{-t^2}\,dt}
 =n+\frac12.
\]
The last equality follows by integrating the derivative of
$t^{2n+1}e^{-t^2}$; its boundary values vanish.
\end{example}

\section{Closed forms and the common invariant domain}
\label{sec:graph-domain}

For a linear operator $S$ on a dense subspace $D\subset H$, the quadratic
form $q(x)=\|Sx\|^2$ is nonnegative.  Its form norm is
$(\|x\|^2+q(x))^{1/2}$.  Such a form is closable if $x_n\to0$ in $H$ and
$q(x_n-x_m)\to0$ imply $q(x_n)\to0$; the implication follows from
closability of $S$.  A closed form has a complete form norm, and the closure
is obtained by completing this norm and identifying the Hilbert limits.
A finite form value permits one operator application.  Repeated products
require control of every algebra element.

For sequence examples, $\ell^2(\mathbb N,w)$ consists of complex sequences
with $\sum_{n\geq1}w_n|x_n|^2<\infty$, for strictly positive weights $w_n$.
Its inner product is $\sum_n w_nx_n\bar y_n$.  Finite sequences, denoted
$c_{00}(\mathbb N)$, are dense by coordinate truncation; $e_n$ denotes the
$n$th coordinate vector.  Unspecified weights are one.
The distinction between a single form and all products already occurs for
a diagonal operator.

\Needspace{5\baselineskip}
\begin{example}[A maximal form domain need not be invariant]
\label{ex:maximal-domain}
On $\ell^2(\mathbb N)$, let $Ne_n=ne_n$, initially on $c_{00}(\mathbb N)$.
Its closure and the closed form $q(v)=\|Nv\|^2$ have domain
\[
 V=\left\{v:\sum_{n\geq1}n^2|v_n|^2<\infty\right\}.
\]
Coordinate convergence proves that the maximal diagonal operator is closed.
Finite coordinate truncations approximate every vector in its graph norm,
so they prove the asserted closure.  The vector $v_n=n^{-2}$ belongs to $V$,
but $Nv\notin V$, since $\sum n^4|v_n|^2=\infty$.
For the polynomial algebra generated by $N$, the common domain is instead
\begin{equation}
 V_\infty=\bigcap_{k\geq0}
 \left\{v:\sum_{n\geq1}n^{2k}|v_n|^2<\infty\right\}.
 \label{eq:all-moments}
\end{equation}
Multiplication by each polynomial preserves $V_\infty$.
Truncations converge in every displayed seminorm.
\end{example}

Return to Theorem~\ref{thm:gns}.  Write $D=D_W$, $H=\mathcal H_W$,
and $A_a=\overline{\pi_W(a)}$ for $a\in\mathfrak A$.
The graph topology on $D$ is generated by $\|\xi\|$ and
$\|\pi_W(a)\xi\|$, for all $a\in\mathfrak A$.
A neighborhood bounds finitely many of these seminorms, so the topology
is locally convex.  A representation is called closed when its domain is
complete for this topology.  Completion can require nets, namely families
indexed by a directed set, when there is no countable family of controlling
seminorms.  A net is Cauchy here when it is Cauchy in every graph seminorm.

\Needspace{5\baselineskip}
\begin{theorem}[The exact common graph domain]
\label{thm:graph-domain}
The completion of $D$ in its graph topology identifies with
\begin{equation}
 D_W^{\mathrm{cl}}=\bigcap_{a\in\mathfrak A}\mathcal D(A_a)
 \subset H.
 \label{eq:common-domain}
\end{equation}
On this domain, the maps $\pi_W^{\mathrm{cl}}(a)=A_a|_{D_W^{\mathrm{cl}}}$
form an invariant $*$-representation.  In particular,
\begin{equation}
 \pi_W^{\mathrm{cl}}(a)\pi_W^{\mathrm{cl}}(b)\xi
 =\pi_W^{\mathrm{cl}}(ab)\xi,
 \qquad
 \langle\pi_W^{\mathrm{cl}}(a)\xi,\eta\rangle
 =\langle\xi,\pi_W^{\mathrm{cl}}(a^*)\eta\rangle.
 \label{eq:closed-products}
\end{equation}
Each action is continuous in the graph topology.
For each $a$, the nonnegative form
\[
 q_a([b])=W(b^*a^*ab)=\|\pi_W(a)[b]\|^2
\]
is closable on $H$.  Its closure is $\overline q_a(\xi)=\|A_a\xi\|^2$
on $\mathcal D(A_a)$.  Thus \eqref{eq:common-domain} is exactly the
intersection of these closed form domains.
Each $\pi_W^{\mathrm{cl}}(a)$ has closure $A_a$.
\end{theorem}

\begin{proof}
A graph Cauchy net $(\xi_i)$ has a Hilbert limit $\xi$ and limits
$\eta_a=\lim_i\pi_W(a)\xi_i$.  Closedness of $A_a$ gives
$\xi\in\mathcal D(A_a)$ and $A_a\xi=\eta_a$.
If $\xi=0$, closability gives $\eta_a=0$ for every $a$.
Hence the graph completion injects into the right side of
\eqref{eq:common-domain}.

To prove the reverse inclusion, fix $\xi$ in that intersection and a finite
set $F\subset\mathfrak A$.  Set $c=\sum_{a\in F}a^*a$.
For $u\in D$, the adjoint identity gives the estimate
\begin{equation}
 \sum_{a\in F}\|\pi_W(a)u\|^2
 =\langle\pi_W(c)u,u\rangle
 \leq\|\pi_W(c)u\|\,\|u\|.
 \label{eq:graph-control}
\end{equation}
Since $\xi\in\mathcal D(A_c)$, choose $u_j\in D$ with
$u_j\to\xi$ and $\pi_W(c)u_j\to A_c\xi$.
Apply \eqref{eq:graph-control} to $u_j-u_k$.
Every $\pi_W(a)u_j$, $a\in F$, is Cauchy.  Closedness identifies its
limit with $A_a\xi$.  Thus $D$ approximates $\xi$ in any prescribed
finite collection of graph seminorms.
Directing these approximations by finite sets and positive tolerances
produces one net converging in all graph seminorms.

Let this net be $(\xi_i)$ and fix $b\in\mathfrak A$.
Then $\pi_W(b)\xi_i\to A_b\xi$ and
$\pi_W(a)\pi_W(b)\xi_i=\pi_W(ab)\xi_i\to A_{ab}\xi$ for every $a$.
Closedness gives $A_b\xi\in D_W^{\mathrm{cl}}$ and
$A_aA_b\xi=A_{ab}\xi$.
Linearity follows by the same approximation.
Approximating both vectors in the algebraic adjoint identity proves
the second equality in \eqref{eq:closed-products}.
Continuity follows from
$\|A_aA_b\xi\|=\|A_{ab}\xi\|$ and $\|A_b\xi\|$ itself.

The form norm for $q_a$ is
$(\|\xi\|^2+\|\pi_W(a)\xi\|^2)^{1/2}$.
Its completion is the graph of $A_a$, which proves the form assertion.
Finally $\pi_W(a)\subseteq\pi_W^{\mathrm{cl}}(a)\subseteq A_a$,
so taking closures proves the last statement.
Any closed $*$-representation extending $\pi_W$ contains the limits of
these graph Cauchy nets and agrees with their multiplier limits.
Thus it extends $\pi_W^{\mathrm{cl}}$, which proves minimality.
\end{proof}

The word ``closed'' refers to completeness in the common graph topology.
It does not assert that each restriction to
$D_W^{\mathrm{cl}}$ is a closed operator or is self-adjoint.
Example~\ref{ex:maximal-domain} separates the common domain from the
maximal domain of one multiplier.

\section{Norm and order in a complete involutive algebra}
\label{sec:cstar-prerequisites}

Bounded target algebras impose estimates which an arbitrary algebraic
representation need not satisfy.  We derive those estimates from the norm
before using them in transfer or unitization.
A Banach space is a complete normed vector space.
A $C^*$-algebra is a complex $*$-algebra complete in a norm satisfying
$\|ab\|\leq\|a\|\|b\|$ and $\|a^*a\|=\|a\|^2$.
The involution is isometric: the latter identity first gives
$\|a\|\leq\|a^*\|$, and application to $a^*$ gives equality.
No representation by bounded operators is assumed in this definition.

\Needspace{5\baselineskip}
\begin{lemma}[Adjoining a unit with its norm]
For\label{lem:cstar-adjoin} a nonunital $C^*$-algebra $A$, the algebra $A^+=A\oplus\C1$ with
\[
 (a+\lambda1)(b+\mu1)=ab+\lambda b+\mu a+\lambda\mu1,
 \qquad (a+\lambda1)^*=a^*+\bar\lambda1
\]
has a complete $C^*$-norm extending the norm of $A$.
If $A\ne\{0\}$, that norm is
\[
 \|a+\lambda1\|_+=\sup_{\|b\|\leq1}\|ab+\lambda b\|.
\]
For $A=\{0\}$ take $A^+=\C$ with the absolute value.
\end{lemma}

\begin{proof}
For $A=\{0\}$ the assertion is the usual $C^*$-norm on $\C$.
Assume henceforth that $A\ne\{0\}$.
Write $L_xb=xb$ for $x\in A^+$ and $b\in A$.
These are bounded maps on the Banach space $A$.
For $a\in A$, the choice $b=a^*/\|a\|$ shows
$\|L_a\|=\|a\|$ when $a\ne0$.
If $L_{a+\lambda1}=0$ with $\lambda\ne0$, then $-a/\lambda$ is a left
identity.  Its adjoint is a right identity; multiplying the two identities
shows that they agree and give a unit in $A$, a contradiction.
Thus $x\mapsto L_x$ is injective and the displayed expression is a norm.

For $b\in A$,
\[
 \|xb\|^2=\|b^*x^*xb\|
 \leq\|b\|\|x^*xb\|\leq\|L_{x^*x}\|\|b\|^2.
\]
Taking suprema gives
$\|L_x\|^2\leq\|L_{x^*x}\|\leq\|L_{x^*}\|\|L_x\|$.
Applying this also to $x^*$ proves equality of the two adjoint norms and
then the $C^*$-identity.

To check completeness directly, one has $|\lambda|\leq\|L_{a+\lambda1}\|$.
Otherwise $L_a=L_{a+\lambda1}-\lambda I$ would be invertible by the
convergent geometric series.  Choose $e\in A$ with $ae=a$.
Then $L_aL_e=L_a$ forces $L_e=I$, again giving a unit in $A$.
It follows that
\[
 \|a+\lambda1\|_+\leq\|a\|+|\lambda|
 \leq3\|a+\lambda1\|_+.
\]
A Cauchy sequence therefore has Cauchy components in $A$ and $\C$, whose
limits supply its limit in $A^+$.
\end{proof}

For an element $h$ of a nonzero unital Banach algebra, its spectrum $\sigma(h)$
is the set of complex $\lambda$ for which $\lambda1-h$ has no two-sided
inverse.  In a nonunital $C^*$-algebra we use the unitization just constructed;
in particular, the spectrum for $A=\{0\}$ is computed in $A^+=\C$.
In every one of these nonzero ambient algebras, $\lambda1$ is invertible
exactly when $\lambda\ne0$: its inverse is $\lambda^{-1}1$, while zero
cannot have an inverse because $1\ne0$.  Thus $\sigma(0)=\{0\}$.
A self-adjoint element is called positive when its spectrum is contained in
$[0,\infty)$.  The next two lemmas show that this is precisely the algebraic
notion of a square $c^*c$.

\Needspace{5\baselineskip}
\begin{lemma}[The calculus of a self-adjoint element]
\label{lem:cstar-calculus}
For a self-adjoint $h$ in a $C^*$-algebra, the spectrum is real and contained
in $[-\|h\|,\|h\|]$.  Every continuous complex function $f$ on this interval
has a well-defined value $f(h)$ obtained by norm limits of polynomials.
The construction preserves sums, products, constants, and conjugation, and
\[
 \|f(h)\|=\max_{t\in\sigma(h)}|f(t)|.
\]
If $A$ is nonunital, $h\in A$, and $f(0)=0$, then $f(h)\in A$.
A continuous function specified on a smaller interval containing the spectrum
may be extended constantly past its endpoints.  Its value is independent
of that extension by the norm identity.
Every positive element has a positive square root in the same algebra.
Every self-adjoint element is a difference $h=h_+-h_-$ of positive elements
with $h_+h_-=0$ and $\|h_\pm\|\leq\|h\|$.
\end{lemma}

\begin{proof}
For $h=0$, take $f(h)=f(0)1$.  We may therefore assume $\|h\|>0$.
We supply the spectral facts used in the polynomial construction.
In a nonzero unital Banach algebra, the series
$(\lambda1-x)^{-1}=\sum_{n\geq0}\lambda^{-n-1}x^n$ converges for
$|\lambda|>\|x\|$; the same series about any inverse proves openness of the
invertible set and holomorphy of the resolvent.
Thus the spectrum is compact.  It is nonempty: an everywhere-defined
resolvent would be an entire bounded Banach-valued function tending to zero
at infinity.  The circle derivative estimate makes such a function constant,
hence zero, contradicting its inverse identity.
The circle formulas of Lemma~\ref{lem:cauchy} hold for Banach-valued
functions by the same uniform Riemann sums and norm estimates; completeness
supplies their limits.

Polynomial factorization shows $\sigma(p(x))=p(\sigma(x))$.
Indeed, the factors of $p(x)-\lambda1$ commute, and a product of commuting
elements is invertible exactly when all its factors are invertible.
Existence of complex roots for a nonconstant polynomial follows from the
same circle estimate: otherwise its reciprocal would be bounded entire and
constant.  Repeated division gives the factorization used here.
If $r(x)=\max\{|\lambda|:\lambda\in\sigma(x)\}$ and $R>r(x)$, deformation
of a circle from radius greater than $\|x\|$ to radius $R$ gives
\[
 x^n=\frac1{2\pi i}\int_{|\lambda|=R}
           \lambda^n(\lambda1-x)^{-1}\,d\lambda.
\]
The norm estimate and the polynomial spectral identity imply
$r(x)=\lim_n\|x^n\|^{1/n}$.
For a normal element $x$, meaning $x^*x=xx^*$, the $C^*$-identity gives
$\|x^2\|=\|x\|^2$.  Iteration at powers $2^n$ therefore gives
$r(x)=\|x\|$.

For self-adjoint $h$, the exponential series gives
$(e^{ith})^*e^{ith}=1$, so $\|e^{ith}\|=1$ for real $t$.
When $\operatorname{Im}\lambda>0$, the norm-convergent integral
\[
 i\int_0^\infty e^{-it(h-\lambda1)}\,dt
\]
is an inverse of $h-\lambda1$, by integration of its derivative.
Conjugation gives an inverse in the lower half-plane.  Hence the spectrum
of $h$ is real.  For any polynomial $p$, the element $p(h)$ is normal;
thus $\|p(h)\|=\max_{\sigma(h)}|p|$.

Continuous functions on a compact real interval admit uniform polynomial
approximation.  An explicit proof on $[0,1]$ uses
$B_nf(t)=\sum_{k=0}^n f(k/n)\binom nk t^k(1-t)^{n-k}$.
The binomial coefficients with these weights have total mass one, mean
$t$, and variance $t(1-t)/n$.  Splitting the sum into $|k/n-t|<\delta$
and its complement bounds the approximation error by
\[
 \sup_{|u-v|<\delta}|f(u)-f(v)|+
       \frac{2\|f\|_\infty}{4n\delta^2}.
\]
Uniform continuity, followed by $n\to\infty$, proves the assertion;
an affine change of variable treats the stated interval.
The polynomial norm identity now defines $f(h)$ by completion and proves
all its displayed properties.  Subtracting the constant term from the
approximating polynomials proves the assertion when $f(0)=0$.

If $h$ is positive, take $f(t)=\sqrt{\max(t,0)}$ on the full interval.
Then $f(h)^2=h$, since the difference vanishes on $\sigma(h)$.
The element $f(h)$ is itself positive: write it as $q(h)^2$ with
$q(t)=(\max(t,0))^{1/4}$, and use the real spectrum of the self-adjoint
$q(h)$ and the polynomial spectral identity.
Taking $f_\pm(t)=\max(\pm t,0)$ gives the final decomposition and norm
bounds by the same argument.
\end{proof}

\Needspace{5\baselineskip}
\begin{lemma}[Positive squares and order bounds]
\label{lem:cstar-order}
The positive elements form a norm-closed cone whose intersection with its
negative is zero.  They are exactly the elements $c^*c$.
Write $a\leq b$ when $b-a$ is positive.  Conjugation by any element preserves
this order, and
\[
 0\leq a\leq b\ \Longrightarrow\ \|a\|\leq\|b\|,
 \qquad c^*c\leq\|c\|^2 1.
\]
\end{lemma}

\begin{proof}
For self-adjoint $h$ and $R>0$, the condition
$\|h-R1\|\leq R$ says exactly $\sigma(h)\subset[0,2R]$.
This criterion and the triangle inequality show that positive elements
are closed under addition: use radii $\|a\|$ and $\|b\|$ for two nonzero
positive elements.  Scalar multiplication by nonnegative reals preserves
positivity.  A norm-convergent sequence of positive elements has a common
bounding radius, so the cone is closed.  If both $h$ and $-h$ are positive,
the spectrum is $\{0\}$ and the normal-element norm identity gives $h=0$.

For every $x$, the nonzero spectra of $x^*x$ and $xx^*$ coincide.
This follows by scaling the identity
$(1-uv)^{-1}=1+u(1-vu)^{-1}v$ whenever the latter inverse exists.
Suppose first that $-x^*x$ is positive.  Writing $x=u+iv$ with self-adjoint
$u,v$ gives
\[
 xx^*=2u^2+2v^2-x^*x\geq0.
\]
The common nonzero spectrum must then be both positive and negative, hence
empty.  Thus $\|x\|^2=\|x^*x\|=0$.
For a general $x$, put $h=x^*x$, take its negative part $h_-$ from
Lemma~\ref{lem:cstar-calculus}, and set $y=xh_-^{1/2}$.
Commutation in the calculus gives $y^*y=hh_-=-h_-^2$.
The preceding argument forces $y=0$, hence $h_-=0$ and $x^*x\geq0$.
Conversely, a positive element is the square of its self-adjoint square
root, proving the characterization.

If $h=d^*d\geq0$, then $c^*hc=(dc)^*(dc)\geq0$.
The spectral bound for $c^*c$ is $[0,\|c\|^2]$, which proves the second
order estimate.  Finally $b\leq\|b\|1$ implies
$a\leq\|b\|1$ when $0\leq a\leq b$.
The spectra of $a$ and $\|b\|1-a$ are nonnegative, so
$\sigma(a)\subset[0,\|b\|]$ and $\|a\|\leq\|b\|$.
\end{proof}

\Needspace{5\baselineskip}
\begin{lemma}[A positive approximate identity]
\label{lem:cstar-approximate}
Every $C^*$-algebra $A$ has a net $(e_i)$ of positive elements with
$\|e_i\|\leq1$ and $e_i a\to a$, $a e_i\to a$ for every $a\in A$.
Such a net is called a positive contractive two-sided approximate identity.
\end{lemma}

\begin{proof}
For a unital algebra take the constant net $1$; the zero algebra is
immediate.  Suppose now that $A$ is nonunital.  For a finite set
$F\subset A$ and $\epsilon>0$, define in the unitization
\[
 t_F=\sum_{a\in F}(a^*a+aa^*),\qquad
 e_{F,\epsilon}=t_F(t_F+\epsilon1)^{-1}.
\]
The calculus shows $e_{F,\epsilon}\in A$ and
$0\leq e_{F,\epsilon}\leq1$, since $t/(t+\epsilon)$ lies in $[0,1]$
for $t\geq0$ and vanishes at zero.
For $a\in F$, both $a^*a$ and $aa^*$ are bounded above by $t_F$.
Writing $e=e_{F,\epsilon}$, Lemma~\ref{lem:cstar-order} gives
\[
 \|a(1-e)\|^2\leq\|(1-e)t_F(1-e)\|
 \leq\sup_{t\geq0}\frac{\epsilon^2 t}{(t+\epsilon)^2}
 =\frac\epsilon4.
\]
Apply the same bound to $a^*$ to obtain $\|(1-e)a\|^2\leq\epsilon/4$.
Direct finite sets by inclusion and positive $\epsilon$ by decreasing size.
Every fixed $a$ then satisfies both limits.
\end{proof}


\Needspace{7\baselineskip}
\section{Dagger transfer and completion}

\begin{theorem}[Dagger transfer]
\label{thm:dagger}
Let $T:\mathfrak A\to\mathfrak B$ be a $*$-homomorphism of complex
$*$-algebras, let $\omega:\mathfrak B\to\C$ be positive, and put
$W=\omega\circ T$.  Then
\begin{equation}
 \langle a,c\rangle_W
 =\langle T(a),T(c)\rangle_\omega,
 \qquad
 \cN_W=T^{-1}(\cN_\omega).
 \label{eq:pullback}
\end{equation}
Consequently
\begin{equation}
 U_0:D_W\longrightarrow D_\omega,
 \qquad [a]\longmapsto[T(a)],
 \label{eq:gns-map}
\end{equation}
is an isometric injection and extends uniquely to an isometry
\begin{equation}
 U:\mathcal H_W\longrightarrow\mathcal H_\omega
 \label{eq:completion-map}
\end{equation}
with closed range.  On the algebraic common domain,
\begin{equation}
 U_0\pi_W(a)=\pi_\omega(T(a))U_0.
 \label{eq:algebraic-intertwining}
\end{equation}

If $\mathfrak B$ is a $C^*$-algebra, both sides of
\eqref{eq:algebraic-intertwining} extend boundedly and
\begin{equation}
 \|\pi_\omega(b)\|\leq\|b\|,
 \qquad
 \|\pi_W(a)\|\leq\|T(a)\|,
 \label{eq:cstar-bounds}
\end{equation}
while
\begin{equation}
 U\pi_W(a)=\pi_\omega(T(a))U
 \quad\text{on }\mathcal H_W.
 \label{eq:bounded-intertwining}
\end{equation}
\end{theorem}

\begin{proof}
For $a,c\in\mathfrak A$,
\[
 W(c^*a)=\omega(T(c)^*T(a)),
\]
which proves positivity, \eqref{eq:pullback}, and the isometry
\eqref{eq:gns-map}.  Every isometry of pre-Hilbert spaces extends
uniquely to their completions; its range is closed because its domain is
complete.  Multiplicativity of $T$ gives
\eqref{eq:algebraic-intertwining}.

Suppose now that $\mathfrak B$ is a $C^*$-algebra.  Lemma~\ref{lem:cstar-order}, with conjugation by $d$, gives
\[
 \omega(d^*b^*bd)\leq\|b\|^2\omega(d^*d),
\]
which proves the first estimate in \eqref{eq:cstar-bounds}.  Taking
$b=T(a)$ and $d=T(c)$ proves the second.  Continuity then extends
\eqref{eq:algebraic-intertwining} to
\eqref{eq:bounded-intertwining}.  No equality between arbitrary
unbounded closure domains is used.
\end{proof}

\Needspace{5\baselineskip}
\begin{proposition}[Transfer of closed graph domains]
Use\label{prop:graph-transfer} the data of Theorem~\ref{thm:dagger}, and set
$K=U\mathcal H_W\subset\mathcal H_\omega$ and $E=U_0D_W$.
On the dense subspace $E\subset K$, define
\[
 \rho(a)=\pi_\omega(T(a))|_E.
\]
Then $\rho$ is a $*$-representation, and
\begin{equation}
 U D_W^{\mathrm{cl}}=D_\rho^{\mathrm{cl}},
 \qquad
 U\pi_W^{\mathrm{cl}}(a)\xi=\rho^{\mathrm{cl}}(a)U\xi.
 \label{eq:graph-transfer}
\end{equation}
Here every closure defining $D_\rho^{\mathrm{cl}}$ is taken in $K$.
For the ambient closures on $\mathcal H_\omega$ one has
\begin{equation}
 U\mathcal D(\overline{\pi_W(a)})
 \subseteq\mathcal D(\overline{\pi_\omega(T(a))}),
 \qquad
 \overline{\pi_\omega(T(a))}U\xi
 =U\overline{\pi_W(a)}\xi.
 \label{eq:ambient-transfer}
\end{equation}
If $T$ is surjective, $U$ is unitary and
$U D_W^{\mathrm{cl}}=D_\omega^{\mathrm{cl}}$.
Without surjectivity this last domain equality need not hold,
even when $U$ is unitary.
\end{proposition}

\begin{proof}
The algebraic intertwining identifies $\rho(a)$ with
$U\pi_W(a)U^{-1}$ on $E$.  As a map onto $K$, $U$ is unitary.
It identifies the individual graph closures and all graph seminorms.
Theorem~\ref{thm:graph-domain} therefore gives \eqref{eq:graph-transfer}.
If $\xi_j\in D_W$ approaches $\xi$ in the graph norm of
$\pi_W(a)$, its images approach $U\xi$ and have multiplier images
approaching $U\overline{\pi_W(a)}\xi$.
This proves \eqref{eq:ambient-transfer}.
For surjective $T$, one has $U_0D_W=D_\omega$ and $T(\mathfrak A)=\mathfrak B$.
Hence the Hilbert spaces and their full graph topologies agree.
\end{proof}

\Needspace{5\baselineskip}
\begin{example}[A Hilbert isometry can miss the full target domain]
Let\label{ex:transfer-domain} $\mathfrak A=c_{00}(\mathbb N)$ and
$\mathfrak B=\C[n]+c_{00}(\mathbb N)$, with pointwise multiplication
and conjugation.  Here $\C[n]$ denotes polynomial sequences on the
positive integers.  Let $T$ be inclusion and set
\[
 \omega(b)=\sum_{n\geq1}e^{-n}b_n,
 \qquad W=\omega|_{\mathfrak A}.
\]
All sums converge, and both functionals are positive and faithful: their
square values are zero only on the zero element, since every weight is positive.
Both GNS completions identify with
$H=\ell^2(\mathbb N,e^{-n})$, because they contain the finite sequences.
Thus $U$ is the identity on $H$.
Every $\mathfrak A$-multiplier is bounded, so $D_W^{\mathrm{cl}}=H$.
The element $n\in\mathfrak B$ acts by an unbounded multiplier.
Finite truncations give cores for every polynomial multiplier.  Therefore
\[
 D_\omega^{\mathrm{cl}}
 =\left\{v:\sum_{n\geq1}e^{-n}n^{2k}|v_n|^2<\infty
               \text{ for all }k\geq0\right\}.
\]
The vector $v_n=e^{n/2}/n$ belongs to $H$, but fails the condition for
$k=1$.  Thus $U D_W^{\mathrm{cl}}\nsubseteq D_\omega^{\mathrm{cl}}$.
Only the restricted algebra $T(\mathfrak A)$ appears in
\eqref{eq:graph-transfer}.
\end{example}

For a field interpretation, algebra products represent repeated operations
on a common domain.  Theorem~\ref{thm:graph-domain} supplies this domain
after positivity has supplied the Hilbert norm.
Proposition~\ref{prop:graph-transfer} preserves it for the represented
algebra.  A vacuum vector, self-adjoint generators, locality, and a positive
energy condition require their own constructions.

\section{Positive unitization for nonunital
\texorpdfstring{$C^*$}{C-star}-algebras}

The unit missing from a convolution algebra cannot be adjoined to an
arbitrary positive functional without an additional bound.  Let $\omega$
be a positive linear functional on a nonunital $*$-algebra $\mathfrak A$.
Write
$\mathfrak A^+=\mathfrak A\oplus\C1$ and
\begin{equation}
 \omega_m(a+\lambda1)=\omega(a)+m\lambda.
 \label{eq:unit-extension}
\end{equation}
Expanding a square in this algebra gives
\[
 \omega_m((a+\lambda1)^*(a+\lambda1))
 =\omega(a^*a)+\lambda\omega(a^*)+\bar\lambda\omega(a)+m|\lambda|^2.
\]
Requiring real nonnegative values first forces $m\geq0$ and Hermitian
symmetry.  For $m>0$, minimization in $\lambda$ gives the inequality below.
For $m=0$, varying $\lambda$ forces $\omega(a)=0$ for every $a$.
Thus the exact algebraic criterion is
\begin{equation}
 \omega_m\text{ is positive}
 \quad\Longleftrightarrow\quad
 \begin{cases}
 m\in\R,\quad m\geq0,\\
 \omega(a^*)=\overline{\omega(a)},\\
 |\omega(a)|^2\leq m\,\omega(a^*a)&(a\in\mathfrak A).
 \end{cases}
 \label{eq:unit-criterion}
\end{equation}

A bounded representation is nondegenerate when
$\overline{\operatorname{span}\{\pi(a)x:a\in\mathfrak A,\ x\in H\}}=H$.
It is cyclic with vector $\Omega$ when
$\overline{\pi(\mathfrak A)\Omega}=H$.

\Needspace{5\baselineskip}
\begin{theorem}[The nonunital $C^*$ case]
Let\label{thm:cstar-unitization} $\mathfrak A$ be a nonunital $C^*$-algebra and let
$\omega:\mathfrak A\to\C$ be positive.  Then $\omega$ is Hermitian and
bounded, and the functional \eqref{eq:unit-extension} on the unitization
constructed in Lemma~\ref{lem:cstar-adjoin} is positive exactly when
\begin{equation}
 m\in\R,\qquad m\geq\|\omega\|.
 \label{eq:minimal-mass}
\end{equation}
Let $(e_i)$ be any positive contractive two-sided approximate identity.
For the minimal extension $m=\|\omega\|$, its GNS vector
$\Omega=[1]$ satisfies
\begin{equation}
 [e_i]\longrightarrow\Omega,
 \label{eq:approx-vector}
\end{equation}
and the restricted representation of $\mathfrak A$ is bounded,
nondegenerate, and cyclic:
\begin{equation}
 \omega(a)=\langle\pi_\omega(a)\Omega,\Omega\rangle,
 \qquad
 \overline{\pi_\omega(\mathfrak A)\Omega}=\mathcal H_\omega.
 \label{eq:cyclic-vector}
\end{equation}
For $m>\|\omega\|$, the vector $[1]$ of that larger extension is not the
limit in \eqref{eq:approx-vector}.
\end{theorem}

\begin{proof}
Every self-adjoint element is the difference of two positive elements,
so positivity makes $\omega$ Hermitian.  If $\omega$ were unbounded, its
values would be unbounded on the positive unit ball: otherwise the
decomposition of an arbitrary element into four positive parts would
bound it.  One could then choose $a_n\geq0$ with
\[
 \|a_n\|\leq2^{-n},\qquad \omega(a_n)=1.
\]
The norm-convergent positive sum $a=\sum_na_n$ would satisfy
$\omega(a)\geq\sum_{n=1}^N\omega(a_n)=N$ for every $N$, a
contradiction.  Hence $\omega$ is bounded.

Lemma~\ref{lem:cstar-approximate} constructs a positive contractive
approximate identity.  Cauchy--Schwarz and
$e_i a\to a$ give
\begin{equation}
 |\omega(a)|^2
 =\lim_i|\omega(e_i a)|^2
 \leq\|\omega\|\,\omega(a^*a).
 \label{eq:hilbert-bound}
\end{equation}
Equations \eqref{eq:unit-criterion} and \eqref{eq:hilbert-bound} prove
positivity for $m=\|\omega\|$; adding $(m-\|\omega\|)$ times the
positive quotient functional $a+\lambda1\mapsto\lambda$, whose value
on a square is $|\lambda|^2$, proves
sufficiency for every larger $m$.  For a positive functional $\eta$ on a
unital $C^*$-algebra, Cauchy--Schwarz and Lemma~\ref{lem:cstar-order} give
\[
 |\eta(a)|^2\leq\eta(1)\eta(a^*a)
             \leq\eta(1)^2\|a\|^2.
\]
Evaluation at $1$ therefore gives $\|\eta\|=\eta(1)$.
Applying this to $\eta=\omega_m$ and restricting yields $\|\omega\|\leq m$.

It remains to identify the vector $[1]$.  If $\|\omega\|>0$, choose
$\|a\|\leq1$ with $|\omega(a)|$ arbitrarily close to $\|\omega\|$.
Cauchy--Schwarz gives
$|\omega(e_i a)|^2\leq\omega(e_i^2)\omega(a^*a)
\leq\omega(e_i^2)\|\omega\|$.
Since $e_i a\to a$, this forces the lower limit of $\omega(e_i^2)$ to be
at least $\|\omega\|$.  The upper bound is $\|\omega\|\|e_i^2\|\leq\|\omega\|$.
Therefore $\omega(e_i^2)\to\|\omega\|$.
Since $0\leq e_i^2\leq e_i\leq1$ in the unitization, also
$\omega(e_i)\to\|\omega\|$.  Therefore
\[
 \|[1]-[e_i]\|^2
 =\|\omega\|-2\omega(e_i)+\omega(e_i^2)\longrightarrow0,
\]
which is \eqref{eq:approx-vector}; the zero functional is immediate.
Moreover $\pi_\omega(a)\Omega=[a]$, proving cyclicity, and
$\pi_\omega(e_i)[a]=[e_i a]\to[a]$, proving nondegeneracy.  The
$C^*$-estimate in \eqref{eq:cstar-bounds} proves boundedness.  If
$m>\|\omega\|$, the same computation instead gives
\[
 \lim_i\|[1]-[e_i]\|^2=m-\|\omega\|>0.\qedhere
\]
\end{proof}

For $\mathfrak A=\{0\}$, the only functional is $\omega=0$, and
$\mathfrak A^+=\C$.  The extension $\omega_m(\lambda)=m\lambda$
is positive exactly when $m\geq0$.
At minimal mass $m=0$, its GNS Hilbert space and $\Omega$ are both zero,
and the constant approximate identity $e_i=0$ gives $[e_i]=\Omega$.
For $m>0$, the extended GNS space is one-dimensional with
$\|\Omega\|^2=m$; the original zero algebra acts as zero and
$[e_i]=0$ cannot converge to $\Omega$.

\Needspace{5\baselineskip}
\begin{example}[Positive unitization can fail algebraically]
Let $\mathfrak A=c_{00}(\mathbb N)$ with pointwise multiplication and
conjugation, and define
\[
 \omega(a)=\sum_{n\geq1}n a_n.
\]
This is a Hermitian positive functional.  If a positive extension
$\omega_m$ existed, Cauchy--Schwarz in the unitization, applied to the
$n$th point mass $e_n$, would give
\[
 n^2=|\omega_m(e_n)|^2
 \leq\omega_m(1)\omega_m(e_n^*e_n)=mn
\]
for every $n$, which is impossible for finite $m$.
\end{example}


\Needspace{5\baselineskip}
\begin{lemma}[The logarithmic size of the gamma term]
\label{lem:gamma-log-size}
For each fixed $\sigma>0$,
\[
 \operatorname{Re}\psi(\sigma+it)=\log|t|+O_\sigma(1)
 \qquad(|t|\geq1).
\]
Consequently $k_\infty(t)=\operatorname{Re}\psi(1/4+it/2)-\log\pi$
differs from $\log|t|$ by a bounded function outside $[-1,1]$.
\end{lemma}

\begin{proof}
For $|t|\geq1$, put
$h_t(x)=(x+\sigma)/((x+\sigma)^2+t^2)$ on $x\geq0$.
Its derivative changes sign at most once, and its maximum is at most
$1/(2|t|)$.  Thus $\int_0^\infty|h_t'(x)|\,dx\leq1/|t|$.
On each unit interval the fundamental theorem of calculus gives
\[
 \left|h_t(n)-\int_n^{n+1}h_t(x)\,dx\right|
 \leq\int_n^{n+1}|h_t'(x)|\,dx.
\]
The real part of the series in \eqref{eq:digamma} therefore equals
\[
 -\gamma+\lim_{N\to\infty}
 \left(\sum_{n=0}^N\frac1{n+1}
 -\frac12\log\frac{(N+1+\sigma)^2+t^2}{\sigma^2+t^2}\right)
 +O(|t|^{-1}).
\]
By the defining limit of $\gamma$, this is
$\tfrac12\log(\sigma^2+t^2)+O(|t|^{-1})$.
The stated bound follows, and replacing $t$ by $t/2$ changes the logarithm
only by $-\log2$.
\end{proof}

\section{The Gaussian obstruction to finite mass}
\label{sec:weil-mass}

For the full Weil functional, failure of positive unitization follows
from its arithmetic local terms.  It does not require a hypothesis about
the zeros.  The relevant functions approach the convolution identity
while remaining inside $\cA_\zeta$:
\begin{equation}
 g_a(u)=\frac{e^{-u^2/(4a)}}{\sqrt{4\pi a}},
 \qquad \widehat g_a(t)=e^{-at^2},
 \qquad g_a^*\star g_a=g_{2a},\quad a>0.
 \label{eq:gaussian-semigroup}
\end{equation}

\Needspace{5\baselineskip}
\begin{theorem}[No finite positive unitization of the Weil functional]
\label{thm:weil-mass}
As $a\downarrow0$, the functional in \eqref{eq:weil-functional} satisfies
\begin{equation}
 L_\zeta(g_a)
 =\frac{a^{-1/2}\log(1/a)}{4\sqrt\pi}+O(a^{-1/2}).
 \label{eq:weil-gaussian-asymptotic}
\end{equation}
In particular $L_\zeta(g_{2a})>0$ for all sufficiently small $a$, and
\begin{equation}
 \frac{|L_\zeta(g_a)|^2}{L_\zeta(g_a^*\star g_a)}
 \sim\frac{\sqrt2}{4\sqrt\pi}\,a^{-1/2}\log(1/a)
 \longrightarrow\infty.
 \label{eq:weil-unitization-divergence}
\end{equation}
There is no positive extension of $L_\zeta$ to $\cA_\zeta^+$ with a
finite value at the adjoined unit.
\end{theorem}

\begin{proof}
Every derivative of $g_a$ is a polynomial times a Gaussian.
Thus $g_a\in\cA_\zeta$.  Completing the square proves
\eqref{eq:gaussian-semigroup}, with the Fourier normalization
\eqref{eq:operations}.
The polar contribution is $2e^{a/4}=O(1)$.

For $0<a\leq1$, split the Gaussian exponent into two equal parts.
Since $\Lambda_{\mathrm{vM}}(n)\leq\log n$, the absolute prime contribution
is bounded by
\begin{equation}
 \frac{e^{-(\log2)^2/(8a)}}{\sqrt{\pi a}}
 \sum_{n\geq2}\frac{\log n}{\sqrt n}
                       e^{-(\log n)^2/8}.
 \label{eq:prime-gaussian-bound}
\end{equation}
The series is finite: its terms are at most $n^{-2}$ for all sufficiently
large $n$.  Hence this contribution is exponentially small relative to
$a^{-1/2}$.

Put $k_\infty(t)=\operatorname{Re}\psi(1/4+it/2)-\log\pi$.
Lemma~\ref{lem:gamma-log-size} shows that
$k_\infty(t)-\log|t|$ is bounded for $|t|\geq1$.
It is integrable on $[-1,1]$, since $k_\infty$ is smooth there.
Therefore
\[
 \int_{\R}e^{-at^2}\bigl(k_\infty(t)-\log|t|\bigr)\,dt
 =O(a^{-1/2}).
\]
Substitution $s=\sqrt a\,t$ gives
\[
 \frac1{2\pi}\int_{\R}e^{-at^2}\log|t|\,dt
 =\frac{a^{-1/2}}{2\pi}
 \left(\frac{\sqrt\pi}{2}\log(1/a)
          +\int_{\R}e^{-s^2}\log|s|\,ds\right).
\]
The remaining integral is finite.  Adding all three local terms proves
\eqref{eq:weil-gaussian-asymptotic}.
Applying that formula at $a$ and $2a$ proves
\eqref{eq:weil-unitization-divergence}.
A positive unitization of mass $m$ would require
$|L_\zeta(g_a)|^2\leq mL_\zeta(g_{2a})$ by
\eqref{eq:unit-criterion}, which contradicts the divergence.
\end{proof}

\Needspace{5\baselineskip}
\begin{theorem}[Positive finite Gaussian subspaces]
\label{thm:gaussian-subspaces}
Fix distinct positive real numbers $\lambda_1,\ldots,\lambda_r$.
There exists $a_0>0$ such that the Weil form is positive definite on
\[
 E_a=\operatorname{span}_{\C}\{g_{a\lambda_1},\ldots,g_{a\lambda_r}\}
 \qquad(0<a<a_0).
\]
The dimension of $E_a$ is $r$.  This assertion is unconditional.
The threshold $a_0$ may depend on the entire finite family.
\end{theorem}

\begin{proof}
The Gram matrix of the displayed functions is the real symmetric matrix
\[
 G(a)_{jk}=L_\zeta(g_{a(\lambda_j+\lambda_k)}).
\]
Apply \eqref{eq:weil-gaussian-asymptotic} at each of the finitely many
scales $a(\lambda_j+\lambda_k)$.  Entry by entry,
\begin{equation}
 \frac{4\sqrt\pi\,\sqrt a}{\log(1/a)}G(a)
 \longrightarrow M,
 \qquad M_{jk}=\frac1{\sqrt{\lambda_j+\lambda_k}}.
 \label{eq:gaussian-gram-limit}
\end{equation}
For a column $c\in\C^r$, let $c^*$ be its conjugate transpose.  Then
\begin{equation}
 c^*Mc=\frac1{\sqrt\pi}\int_{\R}
       \left|\sum_{j=1}^r c_j e^{-\lambda_j t^2}\right|^2dt.
 \label{eq:limiting-gaussian-gram}
\end{equation}
If this integral vanishes, continuity makes its integrand zero everywhere.
Order the $\lambda_j$ increasingly, divide the resulting identity by
$e^{-\lambda_1t^2}$, and let $t\to+\infty$.
This gives $c_1=0$.  Repeating proves $c=0$.
Thus $M$ is positive definite.
Write $\|c\|_2^2=\sum_j|c_j|^2$ and put
\[
 \mu=\inf_{\|c\|_2=1}c^*Mc.
\]
This infimum is attained: from a minimizing sequence, successively
choose subsequences on which each of its finitely many real and imaginary
coordinates converges.  A bounded real sequence has such a subsequence
by bisecting a bounded interval and retaining at each step a half
containing infinitely many terms.  The nested intervals have lengths
tending to zero, and completeness supplies their common limit.
The resulting coordinate limit still has norm one, and continuity of
the finite quadratic expression gives its value $\mu$.
Positive definiteness at that nonzero limit therefore gives $\mu>0$.

Let $A(a)$ denote the matrix on the left of
\eqref{eq:gaussian-gram-limit}, and set
$\delta(a)=\max_{j,k}|A(a)_{jk}-M_{jk}|$.
The entrywise convergence gives $\delta(a)\to0$.
Finite-sum Cauchy--Schwarz gives, for every $c\in\C^r$,
\[
 |c^*(A(a)-M)c|
 \leq\delta(a)\left(\sum_j|c_j|\right)^2
 \leq r\delta(a)\|c\|_2^2.
\]
Choose $0<a<1$ small enough that $r\delta(a)<\mu/2$.
Then $c^*A(a)c\geq(\mu/2)\|c\|_2^2$.
The scalar relating $A(a)$ to $G(a)$ is positive, so $G(a)$ is positive
definite as well.
The same linear-independence argument applied to the Fourier transforms
$e^{-a\lambda_jt^2}$ proves $\dim E_a=r$.
\end{proof}

For $\lambda_1=1$ and $\lambda_2=4$, the limiting matrix is
\[
 M=\begin{pmatrix}1/\sqrt2&1/\sqrt5\\1/\sqrt5&1/\sqrt8\end{pmatrix},
 \qquad\det M=\frac1{20}>0.
\]
This gives a positive two-dimensional family directly from the local terms.
The quantifiers in Theorem~\ref{thm:gaussian-subspaces} fix the finite family
before choosing its scale.  They do not imply positivity on all of
$\cA_\zeta$ at once.

\Needspace{5\baselineskip}
\begin{corollary}[Excluded positive carriers]
\label{cor:weil-carriers}
There is no representation
\[
 L_\zeta(g)=\langle\pi(g)\Omega,\Omega\rangle
\]
on a positive Hilbert space with $\Omega$ of finite norm in a common
invariant $*$-representation domain.
There is also no $*$-homomorphism $T:\cA_\zeta\to\mathfrak B$ to a
$C^*$-algebra with a positive linear functional
$\omega:\mathfrak B\to\C$ such that $L_\zeta=\omega\circ T$.
\end{corollary}

\begin{proof}
In the vector case, the adjoint identity on the invariant domain gives
\[
 |L_\zeta(g)|^2
 \leq\|\Omega\|^2\|\pi(g)\Omega\|^2
 =\|\Omega\|^2L_\zeta(g^*\star g).
\]
This contradicts \eqref{eq:weil-unitization-divergence}.
In the $C^*$ case, positivity implies boundedness and
$|\omega(b)|^2\leq\|\omega\|\omega(b^*b)$ by
Theorem~\ref{thm:cstar-unitization}, or its unital Cauchy--Schwarz proof.
Substitute $b=T(g)$ and use the same divergence.
\end{proof}

This obstruction concerns the functional's representing vector.
It does not prevent bounded GNS multipliers.  Indeed, under RH,
\eqref{eq:zero-norm} and $|\widehat f(t)|\leq\|f\|_{L^1}$ give
\[
 L_\zeta((f\star b)^*\star(f\star b))
 \leq\|f\|_{L^1}^2 L_\zeta(b^*\star b).
\]
Thus every convolution multiplier is bounded under that hypothesis,
and $D_{L_\zeta}^{\mathrm{cl}}=\mathcal H_{L_\zeta}$.
The functional itself has no finite representing vector.
A positive functional on a proper test algebra, or a positive weight
finite on a specified test ideal, remains compatible with this restriction.
A positive weight assigns values in $[0,\infty]$ to positive elements,
additively and homogeneously for nonnegative scalars.  Its finite test
ideal and all products must be specified before a GNS or field
interpretation can be asserted.

\section{Why an indefinite quotient does not give positivity}

The Hermitian part of the construction exists without positivity.
For a complex $*$-algebra $\mathfrak A$ and a Hermitian functional $W$,
meaning $W(a^*)=\overline{W(a)}$, put
\[
 B_W(f,g)=W(g^*f),\qquad
 R_W=\{f:B_W(f,g)=0\text{ for every }g\in\mathfrak A\}.
\]
Then $R_W$ is a left ideal, since
\[
 B_W(af,g)=W(g^*af)=B_W(f,a^*g)=0\quad(f\in R_W).
\]
The quotient $\mathfrak A/R_W$ consequently has a nondegenerate Hermitian
pairing and an algebraic left action satisfying
$B_W(af,g)=B_W(f,a^*g)$.
This applies unconditionally to $(\cA_\zeta,L_\zeta)$ by
\eqref{eq:hermitian}.  It supplies neither a positive norm nor a class of
the missing unit.  Under positivity, Cauchy--Schwarz identifies $R_W$
with $\cN_W$, and Theorem~\ref{thm:gns} applies.

Let $V=\C^2$, $J=\operatorname{diag}(1,-1)$, and
\[
 [x,y]=y^*Jx.
\]
The vectors $e_1+e_2$ and $e_1-e_2$ are isotropic, while their sum is
not.  Thus the isotropic cone $\{x:[x,x]=0\}$ is not a linear subspace
and cannot be divided out.  The radical
$\{x:[x,y]=0\ \forall y\}$ is zero here.  In general, quotienting by the
radical produces only a nondegenerate Hermitian space; it supplies
neither a fundamental symmetry nor a complete Hilbert majorant.
A fundamental symmetry is a linear involution $J$ satisfying
$[Jx,y]=[x,Jy]$ and $[Jx,x]>0$ for $x\ne0$.
A complete Hilbert majorant is a complete Hilbert norm on the same space
for which the original Hermitian pairing is continuous.
In this finite-dimensional example the displayed $J$ gives both choices,
but the radical quotient itself selects neither choice.

Nor is a vector functional for an indefinite adjoint positive.  On
$M_2(\C)$ put $A^\sharp=JA^*J$, let $\Omega=e_1$, and set
$\omega(A)=[A\Omega,\Omega]$.  For
\[
 X=\begin{pmatrix}0&0\\1&0\end{pmatrix}
\]
one has
\begin{equation}
 \omega(X^\sharp X)=[X\Omega,X\Omega]=[e_2,e_2]=-1.
 \label{eq:indefinite-negative}
\end{equation}
Hence an indefinite pairing does not provide the positivity required in
\eqref{eq:weil-criterion}.

\Needspace{7\baselineskip}
\section{The independent arithmetic question}

\begin{proposition}[Formal factorization]
For a linear functional $W$ on a complex $*$-algebra $\mathfrak A$, the
following are equivalent:
\begin{enumerate}
\item $W$ is positive;
\item there are a complex $*$-algebra $\mathfrak B$ and maps
\[
 \omega:\mathfrak B\longrightarrow\C,
 \qquad T:\mathfrak A\longrightarrow\mathfrak B,
\]
where $\omega$ is positive, $T$ is a $*$-homomorphism, and
$W=\omega\circ T$.
\end{enumerate}
\end{proposition}

\begin{proof}
Dagger transfer proves $(\mathrm{ii})\Rightarrow(\mathrm{i})$.  For the converse take
$\mathfrak B=\mathfrak A$, $\omega=W$, and $T$ equal to the identity.
\end{proof}

For $W=L_\zeta$, formal factorization is therefore equivalent to the
positivity in \eqref{eq:weil-criterion}, and hence to RH.  The identity
factorization, or a Hilbert space whose spectral data are prescribed from
the zeta zeros, cannot establish that positivity independently.

\Needspace{5\baselineskip}
\begin{openproblem}[Independent dagger realization of $L_\zeta$]
Construct a complex $*$-algebra $\mathfrak B$, a positive linear
functional $\omega:\mathfrak B\to\C$, and a $*$-homomorphism
$T:\cA_\zeta\to\mathfrak B$ such that
\begin{equation}
 L_\zeta=\omega\circ T.
 \label{eq:factorization}
\end{equation}
Derive the carrier, the map, and the positivity of $\omega$ from the
polar, archimedean, and prime-power structures in
\eqref{eq:weil-functional}, without using RH, the multiset of zeta zeros,
or an assertion equivalent to \eqref{eq:weil-criterion}.  Prove
\eqref{eq:factorization} with all three local terms and the Fourier and
measure conventions above.  Corollary~\ref{cor:weil-carriers} excludes
a positive functional on a $C^*$-algebra as the target.
Any admissible positive functional $\omega$ must fail the finite-mass
unitization bound on the image of $T$.
Theorem~\ref{thm:graph-domain} and
Proposition~\ref{prop:graph-transfer} then specify the common invariant
domain supplied by this factorization.
\end{openproblem}

The resulting representation would still need an additional comparison
to a spectral determinant.  Equality of functionals does not preserve
trace multiplicities.  For example, on $\C$ the positive functional
$W(z)=mz$, with integer $m\geq2$, has a one-dimensional GNS space.
A scalar $\lambda$ consequently has determinant $s-\lambda$ there.
The trace realization $z\mapsto zI_m$ on $\C^m$ has the same functional
and determinant $(s-\lambda)^m$.  Thus a determinant realization must
identify the operator, its domain, and spectral multiplicities separately.


\section{Finite matrix kernels and positive trace}
\label{sec:weil-formal-kernels}

A scalar obtained by closing an operator kernel need not define a
positive pairing.  This distinction can already be seen on finite
matrices inside the formal operator algebra.
For this comparison, use the underlying field $\C$ with the discrete
topology and put
\[
 H_{\mathrm{for}}=\C[[x,y]]/\C.
\]
Enumerate the nonconstant monomials as $e_0,e_1,\ldots$ in increasing
total degree.  The formal topology makes $H_{\mathrm{for}}$ the product
$\C^{\mathbb N}$ of discrete coefficient spaces.  Write
$\epsilon^j(h)=h_j$ for its coefficient functionals.
A continuous endomorphism has a matrix with finitely many nonzero entries
in each row.  Indeed, every output coefficient must depend on finitely
many input coefficients.  Conversely, this condition makes every finite
output projection continuous.  Denote this algebra by $\mathscr E$.
Its operator topology requires agreement of each entire row in a fixed
finite set of output rows.

Put $E_{ij}=e_i\otimes\epsilon^j$ and let $\mathscr F_0$ be their
finite complex span.  These are the matrices with finitely many nonzero
entries in total.  They satisfy
\[
 E_{ij}E_{k\ell}=\delta_{jk}E_{i\ell},\qquad
 E_{ij}^*=E_{ji}.
\]
Conjugating coefficients defines the displayed involution on
$\mathscr F_0$.  This algebra has no unit, since each of its elements
annihilates all but finitely many basis vectors.
It is dense in $\mathscr E$: retaining finitely many whole rows of a
row-finite matrix gives a matrix in $\mathscr F_0$ with those rows unchanged.

The same matrix units act on $D_0=\C^{(\mathbb N)}$, the finite sequences,
inside the Hilbert space $\ell^2(\mathbb N)$ over the usual complex field.
Its inner product is $\langle v,w\rangle=\sum_jv_j\overline{w_j}$.
Every finite matrix preserves $D_0$, and conjugate transpose is its
Hilbert adjoint restricted to $D_0$.  Thus this comparison preserves the
algebra and involution on $\mathscr F_0$.  It does not identify the two
topologies or their completions.

\Needspace{5\baselineskip}
\begin{proposition}[The positive finite-matrix trace]
\label{prop:weil-matrix-trace}
Define $\tau(B)=\sum_i B_{ii}$ on $\mathscr F_0$.
It is a positive linear functional with zero GNS null space.
Its GNS completion is $\ell^2(\mathbb N\times\mathbb N)$, with $[B]$
identified with its coefficient array $(B_{ij})$.
The GNS action is left matrix multiplication, and
\[
 \|\pi_\tau(A)\|=\|A\|_{\ell^2\to\ell^2}
 \qquad(A\in\mathscr F_0).
\]
Nevertheless, $\tau$ has no positive extension of finite mass to
$\mathscr F_0^+=\C I+\mathscr F_0$.
\end{proposition}
\begin{proof}
All sums in the algebraic calculation are finite, and
\[
 \tau(C^*B)=\sum_{i,j}B_{ij}\overline{C_{ij}},\qquad
 \tau(B^*B)=\sum_{i,j}|B_{ij}|^2.
\]
The second expression vanishes exactly when $B=0$.
For a norm-Cauchy sequence of finite arrays, every coefficient has a
limit $b_{ij}$.  A common norm bound bounds every finite sum of
$|b_{ij}|^2$, so the limit array is square-summable.
Pass to the coefficientwise limit in a finite sum of
$|(B_n)_{ij}-(B_m)_{ij}|^2$, then take the supremum over the finite sums.
The Cauchy estimate gives $\|B_n-b\|\to0$.
Conversely, finite rectangular truncations of a square-summable array
converge in this norm.  This identifies the quotient and completion.
Multiplication by $A$ acts separately on each column, so
\[
 \|AB\|^2=\sum_j\|A B_{\bullet j}\|_{\ell^2}^2
       \leq\|A\|_{\ell^2\to\ell^2}^2\sum_j\|B_{\bullet j}\|_{\ell^2}^2.
\]
The operator norm on the right is finite.  The entrywise
Cauchy--Schwarz inequality bounds it by
$(\sum_{i,j}|A_{ij}|^2)^{1/2}$.
Conversely, take a matrix $B$ whose only nonzero column is a finite
vector $v$.  The ratio $\|AB\|/\|B\|$ is $\|Av\|/\|v\|$.
Since $A$ uses only finitely many input coordinates, such vectors
compute its full operator norm.  This proves equality.

For $n\geq1$, let $Q_n=\sum_{i=0}^{n-1}E_{ii}$.
It is a self-adjoint projection with $\tau(Q_n)=n$.
If a positive extension had value $c$ on $I$, then $c$ would be real
and nonnegative, while
\[
 0\leq\tau^+((I-Q_n)^*(I-Q_n))=c-n
\]
for every $n$.  No finite $c$ satisfies these inequalities.
Equivalently,
$|\tau(Q_n)|^2/\tau(Q_n^*Q_n)=n$ is unbounded.
\end{proof}

The map $(\lambda,B)\mapsto\lambda I+B$ identifies the algebraic
unitization with the displayed subalgebra of $\mathscr E$.
It is injective because a scalar multiple of the infinite identity is
a finite matrix only when that scalar is zero.
On this subalgebra, any linear extension of finite-matrix trace has the
form
\[
 \tau_c(\lambda I+B)=c\lambda+\operatorname{tr}(B).
\]
For $A,B\in\mathscr F_0$, its value on a product is
\[
 \tau_c((\lambda I+A)(\mu I+B))
 =c\lambda\mu+\lambda\operatorname{tr}(B)
       +\mu\operatorname{tr}(A)+\sum_{i,j}A_{ij}B_{ji}.
\]
The sum is finite and is unchanged when the two factors are exchanged.
Thus $\tau_c$ is cyclic, but the proposition excludes positivity for every $c$.  In particular, any cyclic trace on
a larger formal differential-operator algebra that restricts to this
formula cannot be positive for an involution restricting to the one above.
This conclusion needs no choice of an involution on the larger algebra.

The finite-dimensional corners $Q_n\mathscr F_0Q_n$ do have positive
traces, and the traces agree under their inclusions.  Their units are
$Q_n$, with masses $n$.  Compatibility of these finite traces supplies
$\tau$ on their union, but supplies no uniform mass bound on a unitization.
The common GNS domain and a vector representing the functional are
therefore separate constructions even in this explicit matrix example.

\Needspace{5\baselineskip}
\begin{proposition}[Formal completion and Hilbert adjoints]
\label{prop:weil-formal-adjoint}
Conjugate transpose on $\mathscr F_0$ does not extend to a continuous
involution on $\mathscr E$ with its operator topology.
The finite-matrix action on $D_0$ also does not extend by the same
coefficient formula to an action of all $\mathscr E$ on
$\ell^2(\mathbb N)$.
\end{proposition}
\begin{proof}
Let $T_n=\sum_{i=0}^nE_{i0}$.
Every fixed row stabilizes, so $T_n$ converges in $\mathscr E$ to
\[
 T(h)=h_0(1,1,1,\ldots).
\]
The matrices $T_n^*=\sum_{i=0}^nE_{0i}$ do not converge in
$\mathscr E$.  Their zeroth rows do not stabilize, and an entire output
row is discrete in the operator topology.  Continuity of a proposed
involution would send the convergent sequence $T_n$ to a convergent
sequence, which is impossible.
For the second assertion, $e_0\in\ell^2$ but
$T(e_0)=(1,1,\ldots)\notin\ell^2$.
\end{proof}

There is also a direct norm obstruction.
The projections $E_{nn}$ tend to zero in the formal operator topology,
but each has Hilbert operator norm one and GNS vector norm one for
$\tau$.  Thus the formal completion, operator-norm completion and
GNS completion are different constructions on the same finite matrices.
For the Weil algebra, Theorem~\ref{thm:graph-domain} specifies a fourth
topology when its multipliers are unbounded: the common graph topology.
None of these completions is selected by cyclicity alone.

\section{Gaussian averages and formal coefficients}
\label{sec:weil-gaussian-coefficients}

The Gaussian in the explicit formula also permits an exact comparison
between a numerical integral and formal polynomial elimination.
It is necessary to specify both the observable algebra and the parameter.
For a positive real number $h$, set
\[
 \mathbb E_h(p)=\frac1{\sqrt{2\pi h}}
       \int_{\R}p(u)e^{-u^2/(2h)}\,du
       =\int_{\R}p(u)g_{h/2}(u)\,du,
 \qquad p\in\C[u].
\]
The polynomial algebra here uses pointwise multiplication and
$p^*(u)=\overline{p(u)}$ for real $u$.
Its involution fixes $u$.  Polynomial growth and Gaussian decay make
every displayed integral absolutely convergent.

\Needspace{5\baselineskip}
\begin{proposition}[Numerical and formal Gaussian elimination]
\label{prop:weil-gaussian-comparison}
For every polynomial $p$ and real $h>0$,
\begin{equation}
 \mathbb E_h(p)
 =\sum_{m\geq0}\frac{h^m}{2^m m!}\,p^{(2m)}(0).
 \label{eq:weil-gaussian-moments}
\end{equation}
The sum is finite.  In particular, $\mathbb E_h$ is positive,
$\mathbb E_h(u)=0$ and $\mathbb E_h(u^2)=h$.

Let $\hbar$ instead be a formal central variable of degree zero.
The same finite derivative formula on each polynomial gives a continuous
$\C[[\hbar]]$-linear map
\[
 \mathsf G:\C[u][[\hbar]]\longrightarrow\C[[\hbar]],\qquad
 \mathsf G\!\left(\sum_{j\geq0}p_j(u)\hbar^j\right)
  =\sum_{j,m\geq0}\frac{p_j^{(2m)}(0)}{2^m m!}\hbar^{j+m}.
\]
Continuity refers to the $\hbar$-adic topologies.
On $\C[u,\hbar]$, evaluating this result at $\hbar=h>0$ agrees
with the numerical Gaussian average of the specialized polynomial.
There is no unital $\C$-algebra homomorphism
$\C[[\hbar]]\to\C$ that sends $\hbar$ to a nonzero number $h$.
\end{proposition}
\begin{proof}
Lemma~\ref{lem:gaussian} gives $\mathbb E_h(1)=1$.
Reflection makes every odd moment zero.
For $n\geq2$, integrate the derivative of
$u^{n-1}e^{-u^2/(2h)}$ over $\R$.
Its values at both ends are zero, and its derivative is integrable.
It follows that
\[
 \mathbb E_h(u^n)=(n-1)h\,\mathbb E_h(u^{n-2}).
\]
Thus
$\mathbb E_h(u^{2m})=(2m)!h^m/(2^m m!)$.
Expansion of $p$ proves \eqref{eq:weil-gaussian-moments}.
The positive integral of $|p|^2$ proves positivity, and the moments
of $u$ and $u^2$ have the stated values.

In the formal formula, the coefficient of $\hbar^N$ uses only
$j,m\geq0$ with $j+m=N$.  There are finitely many such pairs.
Each derivative is a derivative of a polynomial, so this defines every
coefficient without a convergence assumption.
The map preserves every power of $\hbar$, proving continuity.
Reindexing these finite coefficient sums proves
$\C[[\hbar]]$-linearity.  For a polynomial in $u,\hbar$, all sums are
finite, so numerical specialization commutes with the proved moment formula.

For $h\ne0$, the formal series
\[
 (1-\hbar/h)^{-1}=\sum_{n\geq0}(\hbar/h)^n
\]
is an inverse in $\C[[\hbar]]$.
A unital homomorphism sending $\hbar$ to $h$ would send the unit
$1-\hbar/h$ to zero, which is impossible.
At $h=0$, the constant-coefficient homomorphism does exist.
\end{proof}

The comparison at positive $h$ therefore holds on the polynomial
subalgebra.  Extension to a selected class of infinite series requires
an additional summability rule and its compatibility with the operations.
Formal coefficientwise existence alone gives no such rule.

The Gaussian average is not multiplicative:
$\mathbb E_h(u)^2=0$ while $\mathbb E_h(u^2)=h$.
Nor does inserting a Gaussian turn pointwise polynomial multiplication
into convolution.  For $A>0$, the linear map
\[
 J_A:\C[u]\longrightarrow\A_\zeta,\qquad J_A(p)(u)=p(u)g_A(u)
\]
is well-defined, since every derivative is a polynomial times a Gaussian.
But
\[
 J_A(1)\star J_A(1)=g_{2A}\ne g_A=J_A(1),
 \qquad J_A(u)^*=-J_A(u)\ne J_A(u^*).
\]
The first inequality follows already from the values at zero.
The second uses the reflection in the convolution involution.
Thus $J_A$ is neither an algebra homomorphism nor an involution-preserving
map.  Positivity of the numerical Gaussian average cannot pass through
it by Theorem~\ref{thm:dagger}.

\section{The remaining comparison with quantum currents}
\label{sec:weil-current-comparison}

The formal Hamiltonian bracket on $H_{\mathrm{for}}$ is
$[f,g]=f_xg_y-f_yg_x$ modulo constants.
Put $a_f=-f_y\partial_x+f_x\partial_y$ on this quotient.
Changing a representative by a constant changes neither the derivatives
nor the induced operator.  An output through degree $N$ uses only
monomials through degree $N+1$ in each input, because a bracket of
monomials has degree equal to the sum of their degrees minus two.
This proves continuity, including continuity of $f\mapsto a_f$ in the
row topology.  Computing the two derivative coefficients gives
\[
 [a_f,a_g]=-\partial_y[f,g]\,\partial_x
                 +\partial_x[f,g]\,\partial_y=a_{[f,g]}.
\]
Thus these operators give an actual symmetry action on $H_{\mathrm{for}}$.
A numerical positive representation needs further data: a complex
Hilbert space, a common invariant dense domain, and
operators whose restricted adjoints realize the chosen involution.
Proposition~\ref{prop:weil-formal-adjoint} prevents obtaining these data
by completing all finite matrix kernels in the formal operator topology.
Proposition~\ref{prop:weil-gaussian-comparison} also prevents specializing
an arbitrary formal parameter at a positive numerical value.
These are restrictions on the stated carriers and maps.

An algebra generated by products includes the empty product as a unit.
For a positive comparison, one might choose an ordinary degree-zero
subalgebra containing this identity and closed under an involution.
The finite-mass calculation excludes this choice if the positive scalar
functional is to be defined on the entire subalgebra.

\Needspace{5\baselineskip}
\begin{proposition}[An obstruction for unital observable algebras]
\label{prop:weil-unital-current}
There is no unital complex $*$-algebra $\mathfrak B$, positive linear
functional $\omega:\mathfrak B\to\C$, and $*$-homomorphism
$T:\A_\zeta\to\mathfrak B$ such that $L_\zeta=\omega\circ T$.
No continuity hypothesis on $T$ or $\omega$ is needed.
\end{proposition}
\begin{proof}
Suppose these data exist.  Put $m=\omega(1)$, a finite nonnegative real
number.  Extend $T$ to the algebraic unitization by
\[
 T^+(f+\lambda1)=T(f)+\lambda1_{\mathfrak B}.
\]
The unitization product and the multiplicativity of $T$ show that $T^+$
is a unital $*$-homomorphism.
Thus $\omega\circ T^+$ is a positive extension of $L_\zeta$ with mass $m$.
By \eqref{eq:unit-criterion},
\[
 |L_\zeta(g_a)|^2\leq m L_\zeta(g_a^*\star g_a)
                  =m L_\zeta(g_{2a})\qquad(a>0).
\]
Theorem~\ref{thm:weil-mass} makes the ratio on the left to
$L_\zeta(g_{2a})$ unbounded as $a$ decreases to zero.
The denominator is positive for all sufficiently small $a$, by the same
asymptotic.  This contradicts the finite bound $m$.
\end{proof}

A remaining positive arithmetic comparison must therefore choose a
proper nonunital algebra on which the functional is finite.
Its first missing map is an involution-preserving algebra map
\[
 T:\A_\zeta\longrightarrow\mathfrak B
\]
into such a specified algebra of current observables, together with a positive
linear functional $\omega$ on $\mathfrak B$ satisfying
$L_\zeta=\omega\circ T$.
The algebra, involution, domain and positivity of $\omega$ must be
constructed independently of the zeta zero locations.
Theorem~\ref{thm:dagger} would give the isometry
$U:[f]\mapsto[T(f)]$ onto the closed subspace
$K=\overline{[T(\A_\zeta)]}\subset\mathcal H_\omega$.
Proposition~\ref{prop:graph-transfer} identifies the transported common
graph domain with that of the restricted multipliers on
$E=[T(\A_\zeta)]\subset K$, with every closure taken in $K$.
Equality with the full ambient common domain requires a further
hypothesis, such as surjectivity of $T$.
The finite-mass obstruction in Theorem~\ref{thm:weil-mass} still applies:
a finite-norm vector representing this functional, or a positive
functional on a $C^*$-algebra as target, cannot supply this factorization.
A positive functional on a proper nonunital algebra remains a possible
kind of carrier, as the finite-matrix trace demonstrates.  This example
does not supply the map $T$ or the equality with $L_\zeta$.

Any current-field construction must also specify its operator products,
contraction tensors and the scalar trace or weight on those products.
Their convergence with the spatial propagators is an additional analytic
condition.  Neither the strict zero strip nor the equality of the
Gaussian moment formulas constructs this comparison or proves the
positivity of the full Weil form.

\clearpage
\endgroup
